\chapter{METODE PENELITIAN}

\section{Desain Penelitian}
Penelitian ini dilakukan dengan mengadopsi metode \textit{Design and Development} (D\&D) yang dirumuskan oleh \textcite{ellisGuideNoviceResearchers2010}. Metode D\&D didefinisikan sebagai pendekatan sistematis untuk merancang, mengembangkan, dan mengevaluasi suatu artefak yang bisa berupa produk atau sistem sebagai solusi atas permasalahan yang spesifik. Berbeda dengan metode penelitian tradisional yang umumnya berfokus pada pengujian hipotesis atau observasi fenomena, penelitian D\&D menempatkan pengembangan artefak teknologi itu sendiri sebagai kontribusi utama terhadap khazanah pengetahuan. Tujuan utamanya adalah untuk menjembatani kesenjangan antara landasan teoritis dan implementasi praktis melalui penciptaan solusi yang dapat diuji validitas dan efektivitasnya.

Metode ini dipilih karena kesesuaiannya dengan karakteristik penelitian ini yang berfokus pada pengembangan prototipe sistem agentik. Menurut \textcite{ellisGuideNoviceResearchers2010}, prototipe merupakan artefak utama yang dihasilkan dalam penelitian D\&D, terlepas dari sifat artefak yang dikembangkan, baik berupa alat, model, maupun instansiasi baru dari alat atau model yang telah ada. Pelaksanaan penelitian ini mengacu pada enam fase dalam model D\&D, yaitu identifikasi masalah, deskripsi tujuan, perancangan dan pengembangan, pengujian, evaluasi, dan penyampaian hasil. Alur tahapan tersebut divisualisasikan pada Gambar \ref{fig:dnd}.

\begin{figure}[H] 
    \centering
    \includegraphics[width=1\textwidth]{bab3/dnd.png}
    \caption[Tahapan penelitian metode \textit{Design and Development} (D\&D)]{Tahapan penelitian metode \textit{Design and Development} (D\&D) menurut \textcite{ellisGuideNoviceResearchers2010}}
    \label{fig:dnd}
\end{figure}

\section{Identifikasi Masalah}
Identifikasi masalah dilakukan dengan studi literatur mendalam terhadap penelitian-penelitian terkait untuk mengidentifikasi kesenjangan dan peluang penelitian. Metode identifikasi penyakit tanaman konvensional, baik melalui pengamatan visual maupun prosedur laboratorium, menghadapi keterbatasan dalam hal skalabilitas, kecepatan, dan aksesibilitas. Pendekatan berbasis \textit{deep learning} seperti CNN dan YOLO telah terbukti mampu meningkatkan akurasi deteksi, namun masih bergantung pada \textit{supervised learning} dengan data latih berlabel dalam jumlah besar dan kesulitan menangani kelas penyakit baru tanpa pelatihan ulang. Perkembangan VLM membuka peluang untuk mengatasi keterbatasan tersebut melalui kemampuan \textit{zero-shot} dan \textit{multimodal retrieval}, namun pemanfaatannya secara terpadu dalam sistem identifikasi penyakit tanaman masih belum dieksplorasi secara menyeluruh.

Dari sisi sistem agentik, penelitian-penelitian sebelumnya masih didominasi oleh pendekatan \textit{workflow} dengan alur kerja yang bersifat statis, sehingga perilaku sistem kurang adaptif terhadap variasi \textit{query} dan kondisi masukan yang beragam. Belum ada penelitian yang secara komprehensif mengintegrasikan pola ReAct dengan VLM untuk identifikasi penyakit tanaman secara \textit{end-to-end}, yakni dari masukan pengguna hingga pemberian respons, mencakup deteksi objek berbasis \textit{open-vocabulary} dan pengambilan \textit{multimodal}, disertai evaluasi menyeluruh menggunakan pendekatan \textit{LLM-as-a-judge}. Kesenjangan ini menunjukkan perlunya pengembangan artefak berupa sistem agentik yang mengintegrasikan seluruh kapabilitas tersebut sebagai solusi yang lebih adaptif dan fleksibel dalam identifikasi penyakit tanaman.

\section{Deskripsi Tujuan}
Berdasarkan identifikasi masalah yang telah diuraikan, penelitian ini bertujuan untuk mengembangkan sistem agentik \textit{multimodal} yang mampu mengidentifikasi penyakit tanaman secara adaptif. Tujuan tersebut dijabarkan secara lebih rinci dan spesifik sebagai berikut:
\begin{enumerate}[nosep]
    \item Mengembangkan modul deteksi objek tanaman dengan mengombinasikan YOLOv11 untuk kecepatan deteksi \textit{real-time} dan OWLv2 untuk kemampuan \textit{open-vocabulary detection} yang fleksibel tanpa batasan kelas yang ditentukan sebelumnya.
    \item Mengembangkan modul pencarian \textit{multimodal} berbasis model SCOLD untuk identifikasi penyakit tanaman melalui mekanisme \textit{image-text retrieval}, sehingga memungkinkan pengenalan penyakit tanpa ketergantungan pada kelas pelatihan yang kaku.
    \item Mengembangkan modul RAG sebagai mekanisme akses ke basis pengetahuan penyakit tanaman yang relevan, guna mengurangi risiko halusinasi dan meningkatkan keandalan respons agen.
    \item Mengintegrasikan seluruh modul ke dalam sistem agentik berbasis pola ReAct menggunakan ekosistem LangChain, sehingga agen LLM dapat memilih dan menggunakan alat secara adaptif sesuai konteks \textit{query} dan kondisi masukan pengguna.
    \item Mengembangkan antarmuka pengguna berbasis web sebagai media interaksi antara pengguna dan sistem agentik yang telah dibangun.
    \item Mengevaluasi performa sistem agentik yang dikembangkan menggunakan pendekatan \textit{LLM-as-a-judge} dengan kerangka evaluasi yang relevan untuk mengukur kualitas respons sistem secara menyeluruh.
\end{enumerate}

\section{Perancangan dan Pengembangan}
Tahap perancangan dan pengembangan difokuskan pada realisasi teknis dari sistem agentik yang diusulkan. Pada bagian ini, arsitektur sistem dan implementasi komponen-komponen utamanya dijabarkan secara rinci. Pembahasan diawali dengan perancangan sistem agentik yang mengadopsi pola ReAct dan pengembangan agen LLM yang berfungsi sebagai pengendali logika sistem. Selanjutnya, uraian dilanjutkan pada pengembangan modul-modul spesifik yang berperan sebagai bagian dari komponen yang mendukung kapabilitas sistem, meliputi modul deteksi dan modul pengambilan. Akhirnya, pengembangan antarmuka pengguna dipaparkan sebagai komponen interaksi antara pengguna dan sistem yang telah dibangun. Hampir seluruh proses pengembangan dilakukan dalam lingkungan WSL2 (\textit{Windows Subsystem for Linux}) dengan dukungan perangkat yang spesifikasinya disajikan pada Tabel \ref{tab:perangkat_pengembangan}. Spesifikasi perangkat ini menjadi faktor penentu yang memengaruhi pemilihan kebutuhan teknis dalam perancangan dan pengembangan sistem.

{
\setstretch{1.0}
\begin{longtblr}[
    caption = {Spesifikasi perangkat pendukung perancangan dan pengembangan sistem},
    entry = {Spesifikasi perangkat pendukung pengembangan},
    label = {tab:perangkat_pengembangan},
]{
    width = \linewidth,
    colspec = {X[0.5,c] X[1.5,l] X[3,l]},
    hlines, vlines,
    rows = {m},
    rowhead = 1,
    row{1} = {font=\bfseries, halign=c, valign=m},
    rowsep = 4pt,
}
    No. & Komponen & Spesifikasi \\
    1 & Prosesor & AMD Ryzen\texttrademark{} 7 7840HS \\
    2 & Memori & 16 GB DDR5-5600 \\
    3 & Kartu Grafis & NVIDIA\textregistered{} GeForce RTX\texttrademark{} 4060 8 GB \\
\end{longtblr}
}

\textcite{ellisGuideNoviceResearchers2010} menegaskan bahwa tahap perancangan dan pengembangan dalam metode D\&D perlu dilaksanakan secara sistematis dengan mengacu pada metode pengembangan tertentu yang sesuai dengan karakteristik artefak yang dibangun. Metode pengembangan ini diterapkan pada komponen-komponen utama dari sistem agentik, yaitu agen LLM, modul deteksi objek, modul pengambilan informasi, dan antarmuka sistem. Mengingat sistem yang dikembangkan dalam penelitian ini terdiri dari beberapa komponen dengan sifat dan kebutuhan teknis yang berbeda, masing-masing komponen dikembangkan menggunakan metode yang relevan. 

Pengembangan modul deteksi objek mengikuti \textit{Artificial Intelligence Development Life Cycle} (AIDLC) yang dirumuskan oleh \textcite{desilvaArtificialIntelligenceLife2022} karena melibatkan pengembangan model. Sementara itu, pengembangan modul pengambilan informasi dan antarmuka sistem mengikuti \textit{Software Development Life Cycle} (SDLC) dengan model \textit{waterfall} berdasarkan perumusan oleh \textcite{pressmanSoftwareEngineeringPractitioners2020} karena melibatkan rekayasa perangkat lunak konvensional. Adapun komponen agen LLM tidak mengikuti metode pengembangan tertentu karena pengembangannya yang tidak konvensional dan bergantung pada abstraksi tingkat tinggi yang disediakan oleh kerangka kerja yang digunakan.

\subsection{Sistem Agentik}
Sistem agentik dalam penelitian ini dirancang dengan pendekatan agen yang mengadopsi pola ReAct. Implementasi sistem dilakukan menggunakan kerangka kerja LangChain melalui \textit{library} LangGraph \parencite{langchainaiLangChain2025}. Arsitektur sistem yang dikembangkan bersifat \textit{multimodal}, di mana agen mampu menerima dan memproses masukan berupa citra tanaman serta instruksi tekstual dari pengguna secara bersamaan. Kapabilitas eksternal dalam sistem ini dikelompokkan ke dalam dua modul utama berdasarkan kesamaan mekanisme kerjanya, yaitu Modul Deteksi Objek dan Modul Pengambilan Informasi. Modul Deteksi Objek terdiri dari alat deteksi \textit{closed-set} khusus daun menggunakan YOLO dan alat OVD untuk objek tanaman umum menggunakan OWLv2. Sementara itu, Modul Pengambilan Informasi mencakup alat pengambilan citra-teks penyakit tanaman menggunakan SCOLD, alat pengambilan basis data pengetahuan melalui mekanisme RAG, serta alat pencarian internet. Gambar \ref{fig:arsitektur_sistem_agentik} menyajikan arsitektur sistem agentik beserta spesifikasi input dan output pada setiap alat yang terintegrasi.

\begin{figure}[H]
    \centering
    \begin{subfigure}[b]{0.55\textwidth}
        \centering
        \includegraphics[width=\linewidth]{bab3/arsitektur_sistem_agentik.png}
        \caption{}
        \label{fig:arsitektur_sistem_agentik_sub}
    \end{subfigure}
    \vspace{1em}
    \begin{subfigure}[b]{0.65\textwidth}
        \centering
        \includegraphics[width=\linewidth]{bab3/modality_moduL_sistem.png}
        \caption{}
        \label{fig:modality_modul_sistem}
    \end{subfigure}
    \caption[Visualisasi sistem agentik]{Visualisasi sistem agentik: (a) arsitektur sistem agentik secara keseluruhan dengan pola ReAct, (b) \textit{modality} masukan dan keluaran pada setiap modul sistem}
    \label{fig:arsitektur_sistem_agentik}
\end{figure}

Dalam implementasi teknis menggunakan LangGraph, konsep teoritis ReAct yang diajukan oleh \textcite{yaoReActSynergizingReasoning2023} diterjemahkan ke dalam struktur graf berstatus (\textit{stateful graph}). Pada desain aslinya, pola ReAct diimplementasikan secara linier melalui teknik \textit{prompting}, di mana LLM menghasilkan teks terstruktur berisi penalaran (\textit{thought}) dan tindakan (\textit{action}) yang kemudian diurai oleh lingkungan eksternal. Di lain sisi, LangGraph memodelkan agen sebagai sistem yang terdiri dari \textit{node} (simpul) dan \textit{edge} (sisi). 

Struktur ini memetakan siklus penalaran, aksi, dan observasi ke dalam dua komponen utama, yaitu \textit{node} Model yang berisi agen dan \textit{node} Tools yang berisi alat. \textit{node} Model bertanggung jawab untuk menentukan apakah suatu tugas memerlukan penggunaan alat berdasarkan konteks yang diberikan. Jika alat diperlukan, LLM akan menghasilkan permintaan pemanggilan alat (\textit{tool call}) yang valid secara skema untuk dieksekusi secara otomatis oleh \textit{node} Tools. Hasil eksekusi tersebut dikembalikan ke \textit{node} Model sebagai observasi untuk dianalisis lebih lanjut hingga solusi akhir tercapai. 

Status (\textit{state}) agen dikelola sebagai objek yang menyimpan riwayat pesan, yang berfungsi sebagai memori jangka pendek atau \textit{scratchpad}. Alur kerja dimulai dari \textit{node} Model, di mana LLM memproses input dan menentukan langkah selanjutnya. Jika model memutuskan untuk menggunakan alat, alur dialihkan melalui \textit{conditional edge} menuju \textit{node} Tools. Setelah alat dieksekusi, hasilnya dicatat ke dalam state sebagai observasi, dan kendali dikembalikan ke \textit{node} Model untuk mengevaluasi informasi baru tersebut. Proses iteratif ini berlanjut hingga model menilai bahwa informasi yang dikumpulkan telah memadai untuk memberikan jawaban akhir (END). Visualisasi perbandingan antara desain teoritis ReAct dengan implementasi berbasis graf pada LangGraph disajikan pada Gambar \ref{fig:react_langgraph_comparison}.

\begin{figure}[H]
    \centering
    \includegraphics[width=0.8\textwidth]{bab3/reactlanggraph.png}
    \caption[Perbandingan desain teoritis ReAct dan implementasi LangGraph]{Perbandingan antara (a) desain teoritis ReAct oleh \textcite{yaoReActSynergizingReasoning2023} dan (b) implementasi berbasis graf pada LangGraph}
    \label{fig:react_langgraph_comparison}
\end{figure}

\subsection{Agen LLM}
Bagian ini menjelaskan spesifikasi dan konfigurasi agen LLM yang dikembangkan untuk mengimplementasikan pola ReAct pada sistem agentik.

\subsubsection{Konfigurasi}
Agen dalam penelitian ini digerakkan oleh Gemini 3 Flash sebagai LLM utama. Model diakses melalui layanan Gemini API dari \textcite{google2025geminiapi}. Pemilihan model ini didasarkan pada karakteristik model yang menawarkan kecepatan inferensi tinggi dan kemampuan penalaran (\textit{reasoning}) serta penggunaan alat (\textit{tool-use}) yang sangat baik untuk tugas-tugas agentik. Konfigurasi model diatur dengan nilai \textit{temperature} 0 guna memastikan luaran yang bersifat deterministik dan konsisten, sehingga meminimalkan variasi jawaban yang tidak relevan dalam konteks teknis identifikasi penyakit.

Melalui implementasi arsitektur LangGraph, agen memiliki memori jangka pendek (\textit{short-term memory}) yang dikelola melalui objek \textit{State} (status) yang bersifat \textit{stateful}. Status ini berfungsi sebagai buku catatan (\textit{scratchpad}) yang menyimpan seluruh riwayat interaksi dan informasi yang dikumpulkan selama satu sesi tugas berlangsung. Dengan menyimpan status secara terstruktur, agen memiliki konteks terhadap langkah-langkah yang telah diambil sebelumnya, misalnya hasil deteksi objek atau temuan dari pencarian internet, yang kemudian digunakan sebagai dasar untuk menentukan langkah tindakan selanjutnya hingga solusi akhir tercapai. Tabel \ref{tab:state_components} menyajikan rincian variabel status yang dikelola dalam sistem.

{
\setstretch{1.0}
\begin{longtblr}[
    caption = {Komponen status (\textit{state}) pada agen LLM},
    label = {tab:state_components},
]{
    width = \linewidth,
    colspec = {X[1.2,c] X[0.7,c] X[1.2,l]},
    hlines, vlines,
    rows = {m},
    rowsep = 4pt,
    row{1} = {font=\bfseries, halign=c, valign=m},
    column{1} = {font=\bfseries},
    rowhead = 1,
}
    Variabel Status & Tipe Data & Deskripsi \\
    \textit{messages} & \textit{list[AnyMessage]} & Menyimpan riwayat seluruh pesan interaksi (percakapan) antara pengguna dan agen sebagai memori jangka pendek (\textit{scratchpad}). \\
    \textit{images} & \textit{list[dict]} & Menyimpan daftar citra yang diunggah beserta metadata terkait untuk mendukung penanganan beberapa citra sekaligus. \\
    \textit{detections} & \textit{list[dict]} & Menyimpan data hasil deteksi objek yang mencakup label, skor kepercayaan, dan koordinat kotak pembatas (\textit{bounding box}). \\
    \textit{plant\_disease\_classifications} & \textit{list[dict]} & Menyimpan hasil identifikasi penyakit tanaman yang mencakup label prediksi, tingkat kepercayaan, dan metadata pendukung lainnya. \\
    \textit{visualization\_urls} & \textit{list[str]} & Menyimpan daftar tautan publik dari citra visualisasi hasil deteksi. \\
    \textit{current\_image\_url} & \textit{str} & Menyimpan URL (\textit{Uniform Resource Locator}) atau alamat internet dari citra yang sedang menjadi fokus analisis aktif dalam proses iterasi berjalan. \\
    \textit{jump\_to} & \textit{JumpTo} & Atribut kontrol internal untuk mengelola navigasi alur kerja (\textit{routing}) antar simpul (\textit{node}) dalam arsitektur graf LangGraph. \\
\end{longtblr}
}

Untuk memastikan stabilitas, keamanan, dan fungsionalitas yang optimal, agen dilengkapi dengan serangkaian \textit{middleware}. \textit{Middleware} dalam arsitektur ini berfungsi sebagai lapisan perantara yang membungkus pemanggilan model maupun alat (\textit{tools}) untuk pemrosesan tambahan secara otomatis sebelum dan sesudah eksekusi dilakukan. Tabel \ref{tab:middleware_components} merinci komponen \textit{middleware} yang diintegrasikan ke dalam sistem.

{
\setstretch{1.0}
\begin{longtblr}[
    caption = {Komponen \textit{middleware} pada agen LLM},
    label = {tab:middleware_components},
]{
    width = \linewidth,
    colspec = {X[0.8,c] X[1.2,l]},
    hlines, vlines,
    rows = {m},
    rowsep = 4pt,
    row{1} = {font=\bfseries, halign=c, valign=m},
    column{1} = {font=\bfseries},
    rowhead = 1,
}
    Middleware & Deskripsi Fungsi \\
    \textit{get\_system\_prompt} & Menyediakan instruksi dasar, batasan peran, dan panduan penggunaan alat kepada agen. \\
    \textit{handle\_tool\_errors} & Menangkap kegagalan eksekusi pada \textit{tools} dan mengembalikan pesan kesalahan yang informatif agar agen dapat melakukan koreksi secara mandiri. \\
    \textit{model\_retry\_middleware} & Melakukan percobaan ulang otomatis jika terjadi kegagalan koneksi atau batasan kuota (\textit{rate limit}) pada API model. \\
    \textit{tool\_retry\_middleware} & Menangani percobaan ulang jika eksekusi alat eksternal mengalami gangguan teknis sementara. \\
    \textit{model\_fallback\_middleware} & Mekanisme cadangan yang secara otomatis mengalihkan tugas ke model yang lebih tinggi, yakni Gemini 3 Pro apabila model utama gagal merespons. \\
\end{longtblr}
}

\subsubsection{\textit{Prompt Engineering}}
Desain instruksi (\textit{prompt engineering}) pada agen LLM dirancang dengan mengadopsi kerangka kerja CO-STAR (\textit{Context}, \textit{Objective}, \textit{Style}, \textit{Tone}, \textit{Audience}, \textit{Response}) yang didopsi dari \textcite{govtechsingaporeCOSTARPromptEngineering2024}. Kerangka kerja ini digunakan untuk memastikan bahwa model mendapatkan konteks yang lengkap dan terstruktur guna menghasilkan respons yang akurat, relevan, dan konsisten. Gambar \ref{fig:costar_framework} menyajikan visualisasi dari kerangka kerja CO-STAR yang digunakan.

\begin{figure}[H]
    \centering
    \includegraphics[width=0.8\textwidth]{bab3/costar.png}
    \caption{Visualisasi kerangka kerja CO-STAR}
    \label{fig:costar_framework}
\end{figure}

CO-STAR merupakan akronim dari enam komponen utama yang dijelaskan sebagai berikut:
\begin{enumerate}[nosep]
    \item \textit{Context} (Konteks): Memberikan informasi latar belakang mengenai peran agen sebagai asisten identifikasi penyakit tanaman, batasan domain penelitian, serta kapabilitas sistem.
    \item \textit{Objective} (Tujuan): Menetapkan tujuan utama agen untuk menghasilkan identifikasi yang berbasis bukti, jujur terhadap ketidakpastian, dan memberikan panduan yang berlandaskan data.
    \item \textit{Style} (Gaya): Menentukan cara penulisan respons agar mengalir secara alami, eksplisit dalam penyebutan sumber bukti, serta ringkas tanpa pengulangan informasi.
    \item \textit{Tone} (Nada): Mengatur nada bicara agen agar bersifat praktis, tenang, dan profesional tanpa memberikan peringatan yang berlebihan.
    \item \textit{Audience} (Audiens): Mengidentifikasi target pengguna sistem yang mencakup petani, penghobi tanaman, hingga tenaga profesional di bidang pertanian dengan tingkat pemahaman teknis yang bervariasi.
    \item \textit{Response} (Respons): Menentukan struktur luaran yang diharapkan, mulai dari observasi visual, integrasi konteks pengguna, hingga hipotesis identifikasi dan saran tindak lanjut.
\end{enumerate}

Selain komponen standar CO-STAR, desain \textit{prompt} pada sistem ini juga dilengkapi dengan bagian penjelasan alat (\textit{Tools Reference}). Bagian ini berfungsi sebagai dokumentasi teknis yang menjelaskan fungsionalitas, parameter, dan skenario penggunaan dari setiap alat yang terintegrasi, seperti pencarian web, pencarian basis pengetahuan, dan model deteksi visual. Dengan adanya penjelasan alat yang komprehensif, agen dapat memahami cara berinteraksi dengan lingkungan eksternal secara tepat. Seluruh proses perancangan dan iterasi \textit{prompt} dilakukan menggunakan platform LangSmith yang menyediakan antarmuka untuk mengelola, menguji, dan memantau \textit{prompt} secara terpadu.

\subsection{Modul Deteksi Objek}

Modul deteksi objek dirancang sebagai tahap awal dalam alur kerja identifikasi penyakit untuk melakukan lokalisasi terhadap bagian tanaman yang relevan, seperti daun, buah, atau batang. Modul ini memungkinkan agen untuk memperoleh kotak pembatas (\textit{bounding box}) dari objek yang terdeteksi, yang kemudian digunakan oleh modul pengambilan informasi, lebih spesifiknya alat pengambilan citra-teks, untuk melakukan pengambilan berbasis wilayah (\textit{region-based}). Pendekatan ini bertujuan untuk mereduksi derau latar belakang (\textit{background noise}) serta memitigasi gangguan dari objek non-relevan lainnya dalam satu citra, sehingga meningkatkan akurasi pengambilan informasi. Alur kerja dari modul deteksi objek divisualisasikan pada Gambar \ref{fig:plant_detection_workflow}. 

\begin{figure}[H]
\centering
\includegraphics[width=1\textwidth]{bab3/modul_deteksi_objek.png}
\caption{Alur kerja modul deteksi objek}
\label{fig:plant_detection_workflow}
\end{figure}

Implementasi modul ini menggunakan pendekatan hibrida yang menggabungkan dua alat dengan karakteristik yang berbeda, yakni alat deteksi \textit{closed-set} daun dan alat deteksi OVD. Alat deteksi \textit{closed-set} daun mengadaptasi arsitektur YOLOv11 dari \textcite{khanamYOLOv11OverviewKey2024}, sedangkan alat deteksi OVD mengadaptasi model OWLv2 dari \textcite{mindererScalingOpenVocabularyObject2024}. Integrasi kedua model ini bertujuan untuk menciptakan keseimbangan antara efisiensi komputasi pada deteksi objek frekuensi tinggi dan fleksibilitas pada deteksi objek dengan variansi tinggi.

Mengingat modul ini pada dasarnya melibatkan pengembangan model, maka pengembangan modul ini mengiktuti tahapan dalam AIDLC berdasarkan \textcite{desilvaArtificialIntelligenceLife2022}. Tahapan dalam AIDLC meliputi \textit{problem formulation} (formulasi masalah), \textit{data acquisition} (akuisisi data), \textit{data preparation} (persiapan data), \textit{modeling} (pemodelan), \textit{evaluation} (evaluasi), dan \textit{deployment} (penerapan). Setiap tahapan tersebut dijabarkan secara lebih rinci sebagai berikut.

\subsubsection{\textit{Problem Formulation}}
Tahap formulasi masalah bertujuan untuk mengidentifikasi dan mendefinisikan tantangan teknis yang perlu diatasi dalam perancangan modul deteksi objek, sehingga solusi yang dikembangkan dapat menjawab kebutuhan sistem secara tepat. Permasalahan mendasar dalam pengembangan modul deteksi objek untuk identifikasi penyakit tanaman terletak pada variabilitas visual organ tanaman yang menjadi inang gejala penyakit. Walaupun daun merupakan organ yang paling dominan dalam memperlihatkan gejala, identifikasi yang akurat sering kali memerlukan pemeriksaan pada bagian lain seperti buah, batang, atau bunga. Tantangan teknis muncul dalam upaya menyeimbangkan kebutuhan kinerja sistem antara kecepatan (\textit{speed}) dan fleksibilitas (\textit{flexibility}).

Deteksi daun membutuhkan model dengan \textit{throughput} tinggi dan latensi rendah karena merupakan operasi yang paling sering dilakukan oleh sistem. Akan tetapi, model deteksi objek konvensional (\textit{closed-set}) seperti keluarga YOLO, meskipun cepat, terbatas pada kelas yang telah ditentukan saat pelatihan dan kurang adaptif terhadap variasi jenis organ tanaman yang beragam di luar daun. Di sisi lain, model OVD menawarkan fleksibilitas untuk mendeteksi berbagai objek melalui instruksi teks (\textit{text prompt}) tanpa pelatihan ulang, tetapi umumnya memiliki beban komputasi yang jauh lebih berat karena penggunaan dua model sekaligus, yakni \textit{encoder} teks dan citra, yang keluarannya diproses melalui mekanisme fusi informasi untuk dapat bekerja secara \textit{multimodal}. Oleh karena itu, formulasi masalah pada tahap ini difokuskan pada perancangan mekanisme deteksi yang mampu menangani deteksi daun secara efisien sekaligus mengakomodasi deteksi organ tanaman lain secara fleksibel tanpa mengorbankan responsivitas sistem secara keseluruhan.

\subsubsection{\textit{Data Acquisition}}
Tahap akuisisi data mencakup identifikasi dan pengumpulan dataset yang sesuai untuk melatih model deteksi daun pada modul ini. Dataset yang digunakan untuk melatih model deteksi daun diadopsi dari dataset PlantDoc yang disusun oleh \textcite{singhPlantDocDatasetVisual2020}. Dataset ini diakses melalui repositori publik \textcite{singh2019plantdoc}. Pemilihan dataset ini didasarkan pada cakupannya yang luas terhadap variasi spesies tanaman dan kondisi penyakit dalam lingkungan alami, yang sangat relevan untuk aplikasi di dunia nyata. Secara orisinal, dataset ini terdiri dari 2.569 citra yang merepresentasikan 13 spesies tanaman dan terbagi ke dalam 30 kelas kondisi penyakit dan kesehatan. Tampilan repositori sumber dataset ini diperlihatkan pada Gambar \ref{fig:plantdoc_repo}.

\begin{figure}[H]
    \centering
    \begin{subfigure}[b]{1\textwidth}
        \centering
        \includegraphics[width=\textwidth]{bab3/plantdoc_repository.png}
        \caption{}
        \label{fig:plantdoc_repo_view}
        \vspace{0.5cm}
    \end{subfigure}
    \begin{subfigure}[b]{1\textwidth}
        \centering
        \includegraphics[width=\textwidth]{bab3/plantdoc_sampel.png}
        \caption{}
        \label{fig:plantdoc_sample_view}
    \end{subfigure}
    \caption[Dataset PlantDoc: (a) tampilan repositori dan (b) sampel citra]{Dataset PlantDoc: (a) tampilan repositori dan (b) sampel citra}
    \label{fig:plantdoc_repo}
\end{figure}

\subsubsection{\textit{Data Preparation}}
Tahap persiapan data mencakup seluruh proses transformasi dan penyesuaian dataset agar siap digunakan untuk melatih model. Data yang diperoleh dari \textcite{singh2019plantdoc} telah dilengkapi dengan anotasi kotak pembatas (\textit{bounding box}) dan telah diubah ukurannya (\textit{resized}) menjadi resolusi $416 \times 416$ piksel pada repositori sumber. Resolusi ini kompatibel dengan arsitektur masukan standar YOLOv11. Tahap persiapan data dalam penelitian ini difokuskan pada penyatuan seluruh label kelas spesifik yang ada menjadi satu kelas tunggal, yakni "\textit{leaf}" (daun). Langkah ini dilakukan untuk memfokuskan tujuan model pada tugas lokalisasi daun sebagai \textit{Region of Interest} (RoI) utama, terlepas dari jenis penyakitnya. Distribusi data awal dari dataset PlantDoc sebelum penyatuan kelas disajikan pada Gambar \ref{fig:plantdoc_distribution} untuk memberikan gambaran keragaman data yang digunakan. Setelah penyatuan kelas, dataset kemudian dibagi menjadi dua bagian, dengan 90\% data dialokasikan untuk set pelatihan dan 10\% sisanya untuk set pengujian. 
\begin{figure}[H]
    \centering
    \includegraphics[width=1\textwidth]{bab3/plantdoc_distribution.png}
    \caption[Distribusi kelas pada dataset PlantDoc]{Distribusi kelas pada dataset PlantDoc dari \textcite{singhPlantDocDatasetVisual2020}}
    \label{fig:plantdoc_distribution}
\end{figure}

\subsubsection{\textit{Modeling}}
Tahap pemodelan mencakup pemilihan arsitektur model, konfigurasi pelatihan, dan penentuan strategi adaptasi model untuk kedua alat deteksi yang dikembangkan dalam modul ini. Model yang digunakan pada alat deteksi \textit{closed-set} daun dilatih secara \textit{supervised learning} menggunakan pendekatan \textit{transfer learning} pada dataset PlantDoc, sedangkan model yang digunakan pada alat OVD hanya memuat bobot \textit{pre-trained} tanpa pelatihan ulang karena memanfaatkan \textit{foundation model} yang bersifat \textit{open-ended} sehingga tidak memerlukan penyesuaian pada kelas tertentu.

\begin{enumerate}[label=\alph*., listparindent=\parindent, labelindent=0pt]
    \item Alat Deteksi \textit{Closed-Set} Daun

Alat deteksi \textit{closed-set} dalam modul ini diimplementasikan menggunakan arsitektur YOLOv11. Pemilihan YOLOv11 didasarkan pada performanya yang unggul dalam menyeimbangkan antara kecepatan inferensi dan akurasi deteksi dibandingkan versi sebelumnya. Model ini dilatih secara khusus untuk mengenali objek daun di lapangan.

Pemodelan dilakukan dengan mengevaluasi tiga variasi skala model YOLOv11, yaitu YOLOv11n (\textit{nano}), YOLOv11s (\textit{small}), dan YOLOv11m (\textit{medium}). Ketiga variasi model ini dipilih karena kesesuaiannya dengan target lingkungan inferensi berbasis CPU serta kebutuhan sumber daya komputasi pelatihan yang masih dapat diakomodasi oleh perangkat keras yang tersedia dalam penelitian ini. Ketiga varian ini dibandingkan untuk menentukan model yang paling optimal untuk diintegrasikan ke dalam sistem agentik. Seluruh proses pelatihan dilakukan di lingkungan Kaggle Notebooks guna memanfaatkan sumber daya komputasi GPU yang tersedia. Pelatihan dijalankan menggunakan kerangka kerja Ultralytics dengan konfigurasi hiperparameter bawaan, kecuali yang dirinci pada Tabel \ref{tab:yolo_hyperparameters}.

    {
    \setstretch{1.0}
    \begin{longtblr}[
        caption = {Konfigurasi hiperparameter pelatihan YOLOv11},
        entry = {Konfigurasi hiperparameter pelatihan YOLOv11},
        label = {tab:yolo_hyperparameters},
    ]{
        width = \linewidth,
        colspec = {X[l] X[c] X[2,l]},
        hlines, vlines,
        rows = {m},
        rowsep = 3pt,
        rowhead = 1,
        row{1} = {font=\bfseries, halign=c, valign=m},
    }
    Parameter & Nilai & Keterangan \\
    Epoch & 240 & Jumlah iterasi pelatihan penuh. \\
    Input Size & $416 \times 416$ & Resolusi citra masukan sesuai anotasi PlantDoc. \\
    Random Seed & 42 & Menjamin keterulangan hasil eksperimen. \\
    Batch Size & 16 & Jumlah sampel per sub-iterasi. \\
    Patience & 20 & Batas iterasi tanpa peningkatan untuk \textit{early stopping}. \\
    \end{longtblr}
    }

    \item Alat \textit{Open-Vocabulary Object Detection}

Alur OVD (\textit{Open-Vocabulary Object Detection}) diimplementasikan menggunakan model OWLv2 yang dikembangkan oleh Google \parencite{mindererScalingOpenVocabularyObject2024}. Model diakses melalui pustaka Hugging Face Transformers \parencite{huggingface_transformers2026} dengan menggunakan varian google/owlv2-base-patch16. Model ini mengintegrasikan CLIP sebagai \textit{backbone multimodal} dengan arsitektur ViT untuk mendapatkan fitur visual dan model bahasa kausal untuk mendapatkan fitur teks, sehingga mampu melakukan deteksi objek berdasarkan deskripsi tekstual secara \textit{zero-shot}. Gambar \ref{fig:owlv2_info} menyajikan halaman model OWLv2 di Hugging Face beserta visualisasi arsitekturnya.

\noindent
    \begin{minipage}{\linewidth}
        \vspace{0.4cm}
        \captionsetup{type=figure}
        \centering
        \begin{subfigure}[b]{0.8\textwidth}
            \centering
            \includegraphics[width=\textwidth]{bab3/huggingface_owlv2.png}
            \caption{}
            \label{fig:owlv2_huggingface}
            \vspace{0.2cm}
        \end{subfigure}
        \begin{subfigure}[b]{0.8\textwidth}
            \centering
            \includegraphics[width=\textwidth]{bab3/arsitekturalatowlv2.png}
            \caption{}
            \label{fig:arsitektur_owlv2}
        \end{subfigure}
        \caption[Model OWLv2]{Model OWLv2: (a) halaman model di Hugging Face dan (b) arsitektur model}
        \label{fig:owlv2_info}
        \vspace{0.6cm}
    \end{minipage}

\end{enumerate}

\subsubsection{\textit{Evaluation}}
Tahap evaluasi bertujuan mengukur kualitas model yang telah dilatih sebelum diintegrasikan ke dalam sistem. Evaluasi secara khusus dilakukan pada model YOLOv11 yang dilatih. Evaluasi ini dilakukan menggunakan metrik standar deteksi objek yang mencakup ketepatan lokalisasi dan klasifikasi. Metrik yang digunakan meliputi \textit{Precision}, \textit{Recall}, \textit{F1-score}, serta \textit{mean Average Precision} (mAP) pada ambang batas IoU 50\% (mAP50) dan rentang 50--95\% (mAP50-95). Definisi metrik tersebut diuraikan sebagai berikut:

\begin{enumerate}[nosep]
    \item IoU (\textit{Intersection over Union}): Mengukur tingkat tumpang tindih antara kotak prediksi dan \textit{ground truth}. Secara praktis, metrik ini menunjukkan seberapa presisi model dalam melokalisasi posisi objek daun yang sebenarnya sebagaimana ditunjukkan pada Persamaan \ref{eq:iou}.
    \begin{equation}
        IoU = \frac{Area(BB_{pred} \cap BB_{gt})}{Area(BB_{pred} \cup BB_{gt})}
        \label{eq:iou}
    \end{equation}
    Dalam persamaan tersebut, $Area(BB_{pred} \cap BB_{gt})$ merepresentasikan luas area irisan (\textit{intersection}) antara kotak prediksi ($BB_{pred}$) dan kotak \textit{ground truth} ($BB_{gt}$). Sedangkan $Area(BB_{pred} \cup BB_{gt})$ merepresentasikan luas area gabungan (\textit{union}) dari kedua kotak tersebut.
    \daftarpersamaan{Intersection over Union (IoU)}

    \item \textit{Precision}: Rasio prediksi positif yang benar terhadap total prediksi positif. Metrik ini mengindikasikan tingkat kepercayaan terhadap prediksi model, yakni seberapa sering model benar ketika menyatakan mendeteksi daun dan meminimalkan deteksi palsu, sebagaimana dirumuskan pada Persamaan \ref{eq:precision}.
    \begin{equation}
        Precision = \frac{TP}{TP + FP}
        \label{eq:precision}
    \end{equation}
    Dalam persamaan tersebut, $TP$ (\textit{True Positive}) merujuk pada jumlah deteksi daun yang benar, sedangkan $FP$ (\textit{False Positive}) merujuk pada jumlah deteksi yang salah, yaitu ketika model memprediksi adanya daun padahal objek tersebut bukan daun.
    \daftarpersamaan{Precision}

    \item \textit{Recall}: Rasio prediksi positif yang benar terhadap total sampel positif aktual. Metrik ini mengukur sensitivitas model, yakni kemampuannya untuk menemukan seluruh objek daun yang ada dalam citra dan meminimalkan objek yang terlewat, seperti yang dinyatakan dalam Persamaan \ref{eq:recall}.
    \begin{equation}
        Recall = \frac{TP}{TP + FN}
        \label{eq:recall}
    \end{equation}
    Dalam persamaan tersebut, $TP$ (\textit{True Positive}) menyatakan jumlah deteksi daun yang benar, sementara $FN$ (\textit{False Negative}) menyatakan jumlah objek daun yang gagal dideteksi oleh model.
    \daftarpersamaan{Recall}

    \item \textit{F1-score}: Rata-rata harmonis antara \textit{precision} dan \textit{recall}. Metrik ini memberikan gambaran keseimbangan performa model, yang dihitung menggunakan Persamaan \ref{eq:f1}.
    \begin{equation}
        F1 = 2 \cdot \frac{Precision \cdot Recall}{Precision + Recall}
        \label{eq:f1}
    \end{equation}
    \daftarpersamaan{F1-score}

    \item mAP (\textit{mean Average Precision}): Rata-rata \textit{Average Precision} (AP) di seluruh kelas. Metrik ini merupakan standar evaluasi menyeluruh yang merepresentasikan keandalan model pada berbagai tingkat ambang batas deteksi, sebagaimana didefinisikan pada Persamaan \ref{eq:map}.
    \begin{equation}
        mAP = \frac{1}{N} \sum_{i=1}^{N} AP_i
        \label{eq:map}
    \end{equation}
    Dalam persamaan tersebut, $AP_i$ merepresentasikan nilai \textit{Average Precision} untuk kelas ke-$i$, sedangkan $N$ adalah jumlah total kelas yang dievaluasi oleh model.
    \daftarpersamaan{Mean Average Precision (mAP)}
\end{enumerate}

\subsubsection{\textit{Deployment}}
Tahap \textit{deployment} mencakup proses transisi kedua model dari lingkungan pelatihan atau sumber asalnya ke format yang siap diintegrasikan ke dalam sistem agentik. Baik model YOLOv11 maupun OWLv2 diekspor ke format ONNX (\textit{Open Neural Network Exchange}) \parencite{onnx2026}. Konversi ke format ini dilakukan untuk memastikan portabilitas dan kompatibilitas lintas-\textit{framework} dan perangkat. Model YOLOv11 diekspor menggunakan fungsi ekspor bawaan Ultralytics, sedangkan model OWLv2 dikonversi dari bobot Hugging Face Transformers menggunakan Optimum \parencite{huggingface_optimum_onnx2026}. Kedua model selanjutnya dijalankan menggunakan ONNX Runtime. ONNX Runtime merupakan mesin yang berfokus pada performa untuk menjalankan model ONNX dan bertindak sebagai akselerator lintas platform yang menyediakan antarmuka fleksibel untuk integrasi dengan pustaka spesifik perangkat keras \parencite{rothmanRAGDrivenGenerativeAI2024}. Kedua model dioptimalkan dengan \textit{CPU Execution Provider} sesuai dengan target lingkungan inferensi yang berbasis CPU. Akhirnya, berkas model ONNX yang telah dikonversi tersebut dikemas bersama dengan dependensi sistem lainnya ke dalam satu \textit{image} Docker untuk di-\textit{deploy} sebagai bagian dari layanan sistem agentik yang terpadu.

Secara spesifik untuk model YOLOv11, implementasi inferensi melibatkan tahapan prapemrosesan dan pascapemrosesan yang krusial untuk memastikan akurasi deteksi. Tahap prapemrosesan mencakup pengubahan ukuran citra (\textit{resizing}) menjadi resolusi $416 \times 416$ piksel agar sesuai dengan dimensi masukan model. Selain itu, dilakukan normalisasi nilai piksel dengan membaginya dengan 255, sehingga rentang nilai berubah dari $[0, 255]$ menjadi $[0, 1]$. Normalisasi ini penting untuk mempercepat konvergensi perhitungan numerik dan memastikan konsistensi distribusi data masukan terhadap bobot model yang telah dilatih. Setelah model menghasilkan keluaran berupa tensor mentah, tahap pascapemrosesan dilakukan untuk menerjemahkannya menjadi informasi deteksi yang dapat dipahami. Proses ini meliputi penskalaan ulang koordinat kotak pembatas dari dimensi model kembali ke dimensi citra asli, serta penerapan \textit{Non-Maximum Suppression} (NMS). NMS merupakan teknik pascapemrosesan yang berfungsi untuk mengurangi redundansi deteksi dengan menekan kotak pembatas yang saling tumpang tindih, sehingga hanya kotak dengan skor kepercayaan tertinggi yang dipertahankan sebagai representasi terbaik dari objek yang terdeteksi \parencite{redmonYouOnlyLook2016a}.

\subsection{Modul Pengambilan Informasi}
Modul Pengambilan Informasi dirancang sebagai komponen utama sistem untuk mengidentifikasi penyakit tanaman melalui pendekatan pengambilan informasi. Modul ini juga dirancang untuk menyediakan konteks eksternal yang faktual bagi agen dalam menjawab pertanyaan pengguna atau memverifikasi hipotesis identifikasi. Modul ini terdiri dari alat pengambilan citra-teks yang melakukan pengambilan informasi \textit{multimodal} penyakit tanaman dan alat pengambilan basis pengetahuan serta pencarian internet. 

Agen diinstruksikan untuk menggunakan alat pengambilan citra-teks sebagai langkah pertama dalam proses identifikasi penyakit, baik secara langsung atau dengan bantuan hasil deteksi dari modul deteksi objek. Lebih lanjut, agen juga diinstruksikan untuk memanfaatkan alat pengambilan basis pengetahuan serta pencarian internet ketika pengguna meminta panduan manajemen penyakit atau perawatan preventif tanaman. Selain itu, kedua alat ini berfungsi sebagai mekanisme cadangan (\textit{fallback}) faktual terakhir untuk membantu identifikasi apabila alat pengambilan citra-teks penyakit tanaman gagal memberikan hasil yang relevan, misalnya karena penyakit yang diidentifikasi tidak tersedia dalam indeks basis data vektor. 

Alat pengambilan basis pengetahuan dibangun dengan korpus dokumen terkurasi yang berfokus pada penyakit tanaman. Walaupun basis data pengetahuan menawarkan informasi yang lebih terpercaya, cakupannya terbatas pada spesies tanaman yang telah diindeks. Oleh karena itu, pencarian internet melalui Tavily API diintegrasikan sebagai alternatif untuk menangani \textit{query} yang tidak dapat dijawab oleh basis pengetahuan. Agen dirancang untuk memprioritaskan penggunaan pengambilan basis data pengetahuan terlebih dahulu sebelum mengakses informasi dari internet.

Masih mengikuti arahan \textcite{ellisGuideNoviceResearchers2010} untuk menggunakan metode pengembangan yang sesuai dengan jenis artefak, pengembangan modul ini dilakukan dengan mengadopsi model \textit{waterfall} dalam kerangka kerja \textit{Software Development Life Cycle} (SDLC) sesuai rumusan \textcite{pressmanSoftwareEngineeringPractitioners2020}. Pemilihan metode ini didasarkan pada karakteristik pengembangan modul pengambilan informasi yang cenderung bersifat rekayasa perangkat lunak konvensional dengan alur kerja yang linear dan terstruktur. Tahapan proses perangkat lunak (\textit{software process}) menurut \textcite{pressmanSoftwareEngineeringPractitioners2020}, atau yang dirujuk oleh \textcite{sommervilleSoftwareEngineering2016} sebagai SDLC, meliputi \textit{communication} (komunikasi), \textit{planning} (perencanaan), \textit{modeling} (pemodelan), \textit{construction} (konstruksi), dan \textit{deployment} (penerapan). Rincian pelaksanaan setiap tahapan tersebut dijabarkan sebagai berikut.

\begin{figure}[H]
    \centering
    \begin{subfigure}[b]{0.8\textwidth}
        \centering
        \includegraphics[width=\linewidth]{bab3/modul_pengambilan_gambar-teks.png}
        \caption{}
        \label{fig:modul_pengambilan_gambar-teks}
        \vspace{0.2cm}
    \end{subfigure}
    \begin{subfigure}[b]{0.9\textwidth}
        \centering
        \includegraphics[width=\linewidth]{bab3/modul_pengambilan_basis_pengetahuan_dan_internet.png}
        \caption{}
        \label{fig:modul_pengambilan_basis_pengetahuan_dan_internet}
    \end{subfigure}
    \caption[Alur kerja modul pengambilan informasi]{Alur kerja modul pengambilan informasi: (a) pengambilan citra-teks penyakit tanaman dan (b) pengambilan basis pengetahuan dan pencarian interent}
    \label{fig:modul_pengambilan_informasi}
\end{figure}

\subsubsection{\textit{Communication}}
Tahap \textit{communication} bertujuan memahami kebutuhan teknis yang diperlukan sebelum pengembangan dimulai. Pada modul pengambilan informasi terdapat tiga alat, yaitu alat pengambilan citra-teks, alat pengambilan basis pengetahuan, dan alat pencarian internet. Alat pencarian internet pengembangannya cukup diintegrasikan langsung ke dalam sistem dengan memanfaatkan layanan pencarian berbasis API yang dirancang khusus untuk sistem agentik. Dua alat lainnya dikembangkan secara \textit{end-to-end}, yakni dari persiapan data hingga penerapan karena keduanya memerlukan pembangunan \textit{pipeline} RAG yang disesuaikan dengan kebutuhan data domain spesifik.

Kedua alat yang dikembangkan secara \textit{end-to-end} memerlukan basis data vektor yang efisien dan mendukung berbagai teknik pengembangan RAG. Alat pengambilan citra-teks secara khusus membutuhkan basis data vektor yang mampu menyimpan indeks citra dan teks dalam satu koleksi menggunakan representasi vektor terpisah. Untuk keperluan \textit{deployment}, koleksi kedua alat perlu dapat diakses melalui layanan berbasis \textit{cloud}.

Alat pengambilan citra-teks dibangun di atas arsitektur \textit{vanilla} RAG \textit{multimodal}, sehingga memerlukan model \textit{embedding} yang mampu memproses masukan berupa citra dan teks dalam satu ruang vektor yang sama. Alat pengambilan basis pengetahuan menggunakan arsitektur \textit{advanced} RAG dengan teknik \textit{hybrid search} dan \textit{reranking}. Komponen RAG dari alat ini perlu melalui tahap \textit{benchmarking} untuk menentukan konfigurasi terbaik, sehingga memerlukan layanan API eksternal untuk menghasilkan \textit{embedding} vektor padat maupun vektor jarang, layanan model \textit{reranker} berbasis \textit{cross-encoder}, serta kerangka evaluasi untuk mengukur kualitas pengambilan. Selain itu, pengumpulan basis pengetahuan memerlukan proses \textit{web scraping} dengan automasi peramban untuk \textit{crawling} tautan secara sistematis. Ekstraksi konten dari halaman yang dikumpulkan perlu menghasilkan data terstruktur sesuai skema yang telah didefinisikan, sehingga memerlukan LLM yang mampu memahami dan mengurai struktur HTML yang variatif menjadi format terstruktur.

\subsubsection{\textit{Planning}}
Tahap \textit{planning} bertujuan mendefinisikan tumpukan teknologi (\textit{technology stack}) dan sumber daya yang diperlukan sebelum pengembangan dimulai. Berdasarkan kebutuhan yang diidentifikasi pada tahap \textit{communication}, ditetapkan pilihan teknologi untuk masing-masing komponen. Alat pencarian internet dikembangkan menggunakan Tavily Search API \parencite{tavilyaiTavilyaiTavilypython2026}, yaitu mesin pencari yang dioptimasi khusus untuk sistem agentik. Kedua alat yang dikembangkan secara \textit{end-to-end} menggunakan Qdrant sebagai basis data vektor karena efisiensinya serta dukungannya terhadap fitur \textit{named vectors}, yang memungkinkan penyimpanan indeks citra dan teks dalam satu koleksi menggunakan representasi vektor terpisah. Koleksi kedua alat ditempatkan di Qdrant Cloud untuk keperluan \textit{deployment} berbasis \textit{cloud}. Proses \textit{benchmarking} alat pengambilan basis pengetahuan menggunakan Gemini API \parencite{google2025geminiapi} dan OpenAI API \parencite{openai2025api} sebagai penyedia model \textit{dense embedding}, \textit{library} inferensi model untuk \textit{embedding} jenis vektor lainnya, serta Voyage AI \parencite{voyageai2025rerankapi} sebagai penyedia model \textit{reranker} berbasis \textit{cross-encoder}. Evaluasi setiap konfigurasi dilakukan menggunakan ranx \parencite{bassaniRanxBlazingFastPython2022} dengan dukungan LlamaIndex \parencite{runllamaLlamaIndex2025}. Pengumpulan basis pengetahuan dilakukan menggunakan Playwright \parencite{playwright2020} untuk \textit{crawling} dan ScrapeGraphAI \parencite{scrapegraphai2024} untuk ekstraksi konten terstruktur.

Selain penetapan teknologi, kedua alat yang dikembangkan secara \textit{end-to-end} memerlukan sumber dokumen yang perlu diindeks ke dalam basis data vektor. Sumber dan proses persiapan dokumen untuk masing-masing alat diuraikan sebagai berikut.

\begin{enumerate}[label=\alph*., listparindent=\parindent, labelindent=0pt]
    \item Alat Pengambilan Citra-Teks

Dataset PlantWild yang disusun oleh \textcite{weiBenchmarkingInthewildMultimodal2024} digunakan sebagai basis indeks pencarian untuk alat pengambilan citra-teks. Dataset ini mencakup 18.542 citra yang terdistribusi ke dalam 89 kelas kondisi kesehatan dan penyakit pada 13 spesies tanaman, terdiri dari 56 kelas penyakit dan 33 kelas kondisi sehat. Citra dalam dataset ini merepresentasikan berbagai organ tanaman, meliputi daun, buah, dan batang. Sebanyak 3 citra dari setiap kelas diambil secara acak sebagai data uji, menghasilkan 267 citra uji, sedangkan 18.275 citra sisanya digunakan untuk diindeks ke dalam basis data vektor. Selain citra, setiap kelas juga dilengkapi dengan 10 \textit{caption} yang mendeskripsikan karakteristik visualnya. Distribusi kelas dan contoh citra pada dataset PlantWild ditunjukkan pada Gambar~\ref{fig:plantwild_dataset}.

\noindent
\begin{minipage}{\linewidth}
    \vspace{0.4cm}
    \captionsetup{type=figure}
    \centering
    \begin{subfigure}{\linewidth}
        \centering
        \includegraphics[width=\linewidth]{bab3/distribusiplantwild.png}
        \caption{}
        \label{fig:plantwild_distribution}
    \end{subfigure}
    \begin{subfigure}{0.8\linewidth}
        \centering
        \includegraphics[width=0.45\linewidth]{bab3/sampel_plantwild.png}
        \caption{}
        \label{fig:plantwild_samples}
    \end{subfigure}
    \caption[Dataset PlantWild dari \protect\citeauthor{weiBenchmarkingInthewildMultimodal2024}]{(a) distribusi kelas dan (b) sampel data pada dataset PlantWild dari \textcite{weiBenchmarkingInthewildMultimodal2024}}
    \label{fig:plantwild_dataset}
    \vspace{0.6cm}
\end{minipage}

    \item Alat Pengambilan Basis Pengetahuan

Persiapan data korpus untuk alat pengambilan basis pengetahuan mencakup serangkaian tahapan mulai dari pengumpulan data mentah melalui \textit{web scraping}, penyusunan korpus, hingga menghasilkan dataset evaluasi sintetis. Alur kerja persiapan data ini diilustrasikan pada Gambar~\ref{fig:data_preparation_workflow}.

\noindent
\begin{minipage}{\linewidth}
    \vspace{0.4cm}
    \captionsetup{type=figure}
    \centering
    \includegraphics[width=\textwidth]{bab3/flowchart_persiapan_data_rag.png}
    \caption[Alur kerja persiapan data RAG]{Alur kerja persiapan data dari \textit{crawling} hingga pembangkitan dataset QA}
    \label{fig:data_preparation_workflow}
    \vspace{0.6cm}
\end{minipage}

    Data mentah dikumpulkan dari dua sumber utama, yaitu plantvillage.psu.edu \parencite{pennstatePlantVillage2026} dan gardenology.org \parencite{kojianGardenology2026}. plantvillage.psu.edu dipilih karena reputasinya sebagai sumber informasi penyakit tanaman yang kredibel, dikelola oleh Penn State University, sedangkan gardenology.org melengkapi informasi aspek budidaya dan propagasi tanaman yang kurang tercakup di plantvillage.psu.edu. Antarmuka kedua situs ini ditunjukkan pada Gambar~\ref{fig:sumber_data_web}.

    \noindent
    \begin{minipage}{\linewidth}
        \vspace{0.4cm}
        \captionsetup{type=figure}
        \centering
        \begin{subfigure}[b]{0.48\linewidth}
            \centering
            \includegraphics[width=\linewidth]{bab3/plantvillageweb.png}
            \caption{}
            \label{fig:plantvillage_web}
        \end{subfigure}
        \hfill
        \begin{subfigure}[b]{0.48\linewidth}
            \centering
            \includegraphics[width=\linewidth]{bab3/gardenologyweb.png}
            \caption{}
            \label{fig:gardenology_web}
        \end{subfigure}
        \caption[Tampilan antarmuka situs sumber data: (a) plantvillage.psu.edu dan (b) gardenology.org]{Tampilan antarmuka situs sumber data: (a) plantvillage.psu.edu \parencite{pennstatePlantVillage2026} dan (b) gardenology.org \parencite{kojianGardenology2026}}
        \label{fig:sumber_data_web}
        \vspace{0.6cm}
    \end{minipage}

    Pengumpulan data dilakukan menggunakan teknik \textit{web scraping} hibrida yang menggabungkan automasi peramban dan ekstraksi berbasis LLM. Playwright digunakan untuk penemuan dan pengumpulan tautan (\textit{crawling}), sedangkan ScrapeGraphAI digunakan untuk mengurai struktur HTML menjadi format JSON (\textit{JavaScript Object Notation}) terstruktur sesuai skema yang telah didefinisikan. Model Gemini 2.5 Flash dengan temperatur 0 dipilih sebagai LLM pengekstraksi karena jendela konteksnya yang besar dan latensinya yang rendah. Skema data yang digunakan untuk ekstraksi tercantum pada Tabel~\ref{tab:scraping_schema}. Proses ini menghasilkan 222 entri data yang mencakup sekitar 200 spesies tanaman unik.

    {
    \setstretch{1.0}
    \begin{longtblr}[
        caption = {Skema ekstraksi data dari plantvillage.psu.edu dan gardenology.org},
        label = {tab:scraping_schema},
    ]{
        width = \linewidth,
        colspec = {X[1.6,l] X[0.8,c] X[3.2,l]},
        hlines, vlines,
        rows = {m},
        rowsep = 3pt,
        row{1} = {font=\bfseries, halign=c, valign=m},
        column{1} = {font=\bfseries},
        rowhead = 1,
    }
    Atribut & Tipe Data & Deskripsi \\
    \textit{plant} & String & Nama umum tanaman. \\
    \textit{description} & String & Deskripsi umum karakteristik tanaman. \\
    \textit{propagation} & String & Informasi budidaya tanaman. \\
    \textit{pests\_and\_diseases} & List/String & Daftar penyakit beserta gejala, penyebab, dan manajemennya. \\
    \end{longtblr}
    }

    Data JSON hasil \textit{scraping} selanjutnya diproses menjadi korpus dokumen menggunakan LlamaIndex. Setiap entri dipisahkan menjadi dua kategori dokumen, yaitu dokumen budidaya atau propagasi dan dokumen penyakit. Untuk menangani teks panjang, diterapkan \textit{semantic chunking} menggunakan SemanticSplitterNodeParser dari LlamaIndex dengan model \textit{embedding} gemini-embedding-001. Teknik ini menghitung jarak kosinus antar representasi vektor kalimat yang berurutan, lalu hanya melakukan pemisahan dan pemotongan dokumen pada titik-titik dengan perbedaan semantik tertinggi. Parameter \textit{breakpoint percentile threshold} ditetapkan pada nilai 95, sehingga pemisahan hanya terjadi pada 5\% titik dengan divergensi semantik terbesar, yang mengindikasikan pergantian topik secara signifikan. Pendekatan ini memastikan setiap potongan dokumen memiliki koherensi tematik yang terjaga. Proses ini menghasilkan total 2.469 potongan dokumen (\textit{chunks}) yang siap diindeks, dengan contoh sampel ditunjukkan pada Gambar~\ref{fig:sampel_korpus}.

    \noindent
    \begin{minipage}{\linewidth}
        \vspace{0.4cm}
        \captionsetup{type=figure}
        \centering
        \includegraphics[width=\linewidth]{bab3/sampelkorpus.png}
        \caption{Sampel korpus dokumen hasil \textit{semantic chunking}}
        \label{fig:sampel_korpus}
        \vspace{0.6cm}
    \end{minipage}

    Untuk keperluan \textit{benchmarking} konfigurasi RAG, dibangkitkan pula dataset \textit{Question-Answering} (QA) sintetis sebanyak 600 sampel dari korpus yang telah disusun, menggunakan Gemini 2.5 Flash dengan temperatur 0. Proses menghasilkan dataset mengadopsi metodologi AutoRAG dari \textcite{kimAutoRAGAutomatedFramework2024}. Dataset mencakup tiga jenis pertanyaan dengan tingkat kompleksitas berbeda yang dirangkum pada Tabel~\ref{tab:qa_types}.

    {
    \setstretch{1.0}
    \begin{longtblr}[
        caption = {Jenis pertanyaan pada dataset evaluasi QA sintetis},
        entry = {Jenis pertanyaan pada dataset QA sintetis},
        label = {tab:qa_types},
    ]{
        width = \linewidth,
        colspec = {X[1.4,c] X[2,l] X[2.6,l] X[0.9,c]},
        hlines, vlines,
        rows = {m},
        rowsep = 3pt,
        row{1} = {font=\bfseries, halign=c, valign=m},
        rowhead = 1,
    }
    Tipe Pertanyaan & Definisi & Contoh & Proporsi \\
    \textit{Factoid Single Hop} & Pertanyaan yang mencari fakta singkat dan spesifik yang dapat diverifikasi. & "Apa patogen penyebab penyakit hawar daun pada tomat?" & 35\% \\
    \textit{Concept Completion} & Pertanyaan yang menanyakan esensi atau definisi dari suatu konsep. & "Jelaskan karakteristik dari gejala klorosis." & 35\% \\
    \textit{Two-Hop Incremental} & Pertanyaan yang membutuhkan inferensi dari dua segmen informasi berbeda. & "Bagaimana cara menangani penyakit yang disebabkan oleh jamur Fusarium pada tanaman pisang?" & 30\% \\
    \end{longtblr}
}

\end{enumerate}

\subsubsection{\textit{Modeling}}

Tahap \textit{modeling} dalam SDLC menurut \textcite{pressmanSoftwareEngineeringPractitioners2020} berfokus pada analisis dan perancangan dari perangkat lunak yang dikembangkan. Perlu dicatat bahwa penggunaan istilah \textit{modeling} di sini berbeda konteksnya dengan tahap \textit{modeling} pada kerangka AIDLC \parencite{desilvaArtificialIntelligenceLife2022} yang diterapkan pada modul deteksi objek, di mana tahapan tersebut secara spesifik merujuk pada pengembangan model \textit{machine learning}. Pada modul pengambilan informasi, tahap ini mencakup perancangan struktur indeks basis data vektor dan teknik pengambilan untuk masing-masing alat. Rincian perancangan setiap alat diuraikan sebagai berikut.

\begin{enumerate}[label=\alph*., listparindent=\parindent, labelindent=0pt]
    \item Alat Pengambilan Citra-Teks

    Alat ini dibangun menggunakan arsitektur \textit{vanilla} RAG \textit{multimodal} dengan memanfaatkan model \textit{embedding} yang mampu memproses representasi numerik teks dan citra dalam satu ruang vektor. Model yang digunakan adalah SCOLD dari \textcite{huggingfaceSCOLD} yang dikembangkan oleh \textcite{quocVisionLanguageFoundationModel2025}. Model ini mengintegrasikan arsitektur RoBERTa dari \textcite{liuRoBERTaRobustlyOptimized2019} sebagai \textit{encoder} teks dan arsitektur Swin-T dari \textcite{liuSwinTransformerHierarchical2021} sebagai \textit{encoder} citra. Tampilan laman model tersebut di repositori Hugging Face diperlihatkan pada Gambar \ref{fig:huggingface_scold}.

    \noindent
    \begin{minipage}{\linewidth}
        \vspace{0.4cm}
        \captionsetup{type=figure}
        \centering
        \includegraphics[width=0.8\textwidth]{bab3/huggingface_scold.png}
        \caption[Tampilan model SCOLD di Hugging Face]{Tampilan model SCOLD di repositori Hugging Face \parencite{huggingfaceSCOLD}}
        \label{fig:huggingface_scold}
        \vspace{0.6cm}
    \end{minipage}

    Pengembangan alat ini mencakup dua aspek utama, yaitu mekanisme pengindeksan data penyakit tanaman dan teknik pengambilan informasi \textit{multimodal}.

    \begin{enumerate}[label=\arabic*., listparindent=\parindent, labelindent=0pt]
        \item Pengindeksan

        Untuk pengindeksan, setiap entitas penyakit tanaman disimpan sebagai \textit{point} dalam koleksi Qdrant dengan dua vektor padat bernama (\textit{named dense vectors}), yaitu vektor teks dan vektor citra. Citra referensi asli disimpan secara terpisah pada layanan penyimpanan objek Cloudflare R2 \parencite{cloudflare_r2} untuk efisiensi, sementara URL publiknya disimpan dalam \textit{payload}. Skema \textit{point} beserta struktur \textit{payload} yang diimplementasikan disajikan pada Tabel~\ref{tab:qdrant_schema}.

        {
        \setstretch{1.0}
        \begin{longtblr}[
            caption = {Skema \textit{point} (vektor dan \textit{payload}) pada indeks Qdrant alat pengambilan citra-teks},
            label = {tab:qdrant_schema},
        ]{
            width = \linewidth,
            colspec = {X[1.5,l] X[1.2,l] X[0.8,l] X[3,l]},
            hlines, vlines,
            rows = {m},
            rowsep = 3pt,
            row{1} = {font=\bfseries, halign=c, valign=m},
            column{1} = {font=\bfseries},
            rowhead = 1,
        }
        Komponen & Kunci & Tipe & Deskripsi \\
        \SetCell[r=2]{l} Vektor & \textit{text} & 512 & Representasi semantik deskripsi gejala (RoBERTa). \\
         & \textit{image} & 512 & Representasi fitur visual citra tanaman (Swin-T). \\
        \SetCell[r=6]{l} Payload & \textit{plant\_name} & String & Nama spesies tanaman. Diindeks secara \textit{full-text} untuk filtrasi. \\
         & \textit{label} & String & Nama spesifik penyakit. Diindeks secara \textit{full-text}. \\
         & \textit{caption} & String & Deskripsi tekstual lengkap mengenai gejala. \\
         & \textit{image\_url} & String & Tautan publik citra. \\
         & \textit{source\_id} & Integer & ID referensi internal. \\
         & \textit{filename} & String & Nama fail asli citra. \\
        \end{longtblr}
        }

        Fitur penting dalam implementasi ini adalah penerapan indeks teks penuh (\textit{full-text index}) yang disediakan oleh Qdrant pada kolom \textit{payload} \textit{plant\_name} dan \textit{label}. Berbeda dengan pencocokan kata kunci eksak, indeks teks penuh memungkinkan agen melakukan penyaringan (\textit{filtering}) yang lebih fleksibel berbasis token kata dan tidak peka terhadap huruf besar-kecil. Fitur ini secara signifikan meningkatkan efisiensi pengambilan dengan membatasi ruang pencarian vektor hanya pada subset data yang relevan dengan konteks tanaman yang sedang dianalisis, misalnya dengan menerapkan filter \textit{plant\_name} untuk hanya menelusuri penyakit pada spesies tertentu. Visualisasi mekanisme ini disajikan pada Gambar~\ref{fig:payload_filtering}.

        \noindent
        \begin{minipage}{\linewidth}
            \vspace{0.4cm}
            \captionsetup{type=figure}
            \centering
            \includegraphics[width=1\textwidth]{bab3/fulltextindexpayloadfield.png}
            \caption{Visualisasi mekanisme penyaringan \textit{payload} berbasis \textit{full-text index}}
            \label{fig:payload_filtering}
            \vspace{0.6cm}
        \end{minipage}

        \item Teknik Pengambilan

        Dalam tahap pengambilan, agen dapat memilih salah satu dari empat teknik pengambilan \textit{multimodal} sesuai ketersediaan masukan dan tujuan pengambilan. Fleksibilitas ini dimungkinkan oleh arsitektur SCOLD yang menyelaraskan representasi fitur teks dan citra ke dalam satu ruang laten yang sama, sehingga operasi pengambilan lintas \textit{modal} (\textit{cross-modal}) dapat dilakukan. Keempat teknik tersebut diilustrasikan pada Gambar~\ref{fig:multimodal_search_strategies} dan diuraikan sebagai berikut.

        \noindent
        \begin{minipage}{\linewidth}
            \vspace{0.64cm}
            \captionsetup{type=figure}
            \centering
            \begin{subfigure}[b]{0.45\textwidth}
                \centering
                \includegraphics[width=\textwidth]{bab3/strategipencarianteksteks.png}
                \caption{}
                \label{fig:strat_t2t}
            \end{subfigure}
            \hfill
            \begin{subfigure}[b]{0.45\textwidth}
                \centering
                \includegraphics[width=\textwidth]{bab3/strategipencariantekksgambar.png}
                \caption{}
                \label{fig:strat_t2i}
            \end{subfigure}
    
            \vspace{0.5cm}
    
            \begin{subfigure}[b]{0.45\textwidth}
                \centering
                \includegraphics[width=\textwidth]{bab3/strategipencariangambargambar.png}
                \caption{}
                \label{fig:strat_i2i}
            \end{subfigure}
            \hfill
            \begin{subfigure}[b]{0.45\textwidth}
                \centering
                \includegraphics[width=\textwidth]{bab3/strategipencariangambarteks.png}
                \caption{}
                \label{fig:strat_i2t}
            \end{subfigure}
            \caption[Empat teknik pengambilan \textit{multimodal} pada alat pengambilan citra-teks]{Empat teknik pengambilan \textit{multimodal}: (a) \textit{Text-to-Text}, (b) \textit{Text-to-Image}, (c) \textit{Image-to-Image}, dan (d) \textit{Image-to-Text}}
            \label{fig:multimodal_search_strategies}
            \vspace{0.6cm}
        \end{minipage}
    
        \begin{enumerate}[label=(\alph*), listparindent=\parindent, labelindent=0pt]
            \item \textit{Text-to-Text} \\
            Teknik ini digunakan ketika pengguna memberikan deskripsi gejala berupa teks. Masukan diproses oleh \textit{text encoder} RoBERTa untuk menghasilkan vektor teks, kemudian dicocokkan dengan koleksi vektor teks dalam indeks Qdrant yang merepresentasikan \textit{caption} penyakit.
    
            \item \textit{Text-to-Image} \\
            Teknik ini memungkinkan pengambilan citra referensi penyakit dari masukan teks. Vektor teks hasil enkripsi digunakan untuk mencari tetangga terdekat pada ruang vektor citra, yang dimungkinkan karena SCOLD dilatih untuk meminimalkan jarak antara pasangan teks dan citra yang relevan dalam \textit{embedding space}.
    
            \item \textit{Image-to-Image} \\
            Teknik ini merupakan metode standar untuk identifikasi berbasis citra. Citra masukan dari pengguna diproses oleh \textit{image encoder} Swin-Transformer menjadi vektor citra, kemudian dibandingkan dengan vektor citra dalam basis data untuk menemukan penyakit dengan karakteristik visual paling serupa.
    
            \item \textit{Image-to-Text} \\
            Teknik ini membalikkan pendekatan lintas \textit{modal}, di mana vektor citra masukan digunakan untuk mencari kecocokan pada ruang vektor teks. Teknik ini berguna untuk memperoleh deskripsi tekstual atau label penyakit yang paling relevan secara langsung dari fitur visual.
        \end{enumerate}
    \end{enumerate}

    \item Alat Pengambilan Basis Pengetahuan

    Alat ini dikembangkan menggunakan arsitektur \textit{Advanced} RAG yang menerapkan teknik pencarian hibrida dan \textit{reranking}. Teknik pencarian hibrida diimplementasikan dengan memanfaatkan model yang dapat menghasilkan representasi vektor padat  serta vektor jarang, di mana hasil dari kedua representasi tersebut difusikan. Selanjutnya, model \textit{reranker} digunakan untuk melakukan proses pengurutan ulang terhadap dokumen hasil pencarian awal. Perancangan alat ini mencakup dua aspek utama, yaitu mekanisme pengindeksan dokumen basis pengetahuan dan teknik pengambilan informasi menggunakan pendekatan RAG tingkat lanjut. 

    \begin{enumerate}[label=\arabic*., listparindent=\parindent, labelindent=0pt]
        \item Pengindeksan

        Pengindeksan diawali dengan mempersiapkan representasi numerik pada indeks untuk keperluan pengujian dan evaluasi konfigurasi pengambilan yang optimal. Model representasi vektor padat yang disiapkan mencakup gemini-embedding-001 \textcite{googleGeminiEmbedding0012025} dan text-embedding-3-large dari \textcite{openaiTextEmbedding3Large2024}. Selain itu, representasi vektor jarang juga disiapkan menggunakan BM25 yang implementasinya disediakan secara internal oleh Qdrant dan SPLADE \parencite{formalSPLADESparseLexical2021} yang dijalankan melalui FastEmbed \parencite{qdrant_fastembed_2023}. Terkait model \textit{reranking} yang akan diuji, yakni jina-colbert-v2 \parencite{jhaJinaColBERTv2GeneralPurposeMultilingual2024} dan rerank-2.5 \parencite{voyageai2025rerank25}, hanya jina-colbert-v2 yang memerlukan pengindeksan vektor terlebih dahulu. Hal ini dikarenakan arsitektur ColBERT menghasilkan representasi multi-vektor untuk setiap dokumen yang perlu disimpan guna memungkinkan mekanisme \textit{late interaction} yang efisien. Sebaliknya, rerank-2.5 yang berbasis \textit{cross-encoder} memproses pasangan \textit{query} dan dokumen secara bersamaan pada saat inferensi, sehingga tidak memerlukan penyimpanan representasi vektor statis sebelumnya. 

        Setiap dokumen korpus disimpan sebagai \textit{point} dalam koleksi Qdrant dengan empat vektor bernama sekaligus, yang memungkinkan evaluasi berbagai konfigurasi pengambilan dalam satu koleksi terpadu. Setiap dokumen korpus diubah menjadi \textit{point} yang memiliki empat komponen vektor, yaitu dua vektor padat dan dua vektor jarang. Vektor padat dihasilkan oleh model gemini-embedding-001 dan text-embedding-3-large, masing-masing berdimensi 3.072. Sementara itu, vektor jarang dihasilkan menggunakan SPLADE dan algoritma BM25 yang terintegrasi secara internal di Qdrant. Skema representasi data pada indeks ini disajikan pada Tabel~\ref{tab:knowledge_base_schema}.

        {
        \setstretch{1.0}
        \begin{longtblr}[
            caption = {Skema \textit{point} dengan vektor jamak dan metadata pada indeks basis pengetahuan},
            label = {tab:knowledge_base_schema},
        ]{
            width = \linewidth,
            colspec = {X[1.4,l] X[2.45,c] X[0.9,c] X[2.4,l]},
            hlines, vlines,
            rows = {m},
            rowsep = 3pt,
            row{1} = {font=\bfseries, halign=c, valign=m},
            column{1} = {font=\bfseries},
            rowhead = 1,
        }
        Komponen & Kunci & Tipe & Deskripsi \\
        \SetCell[r=2]{l} Vektor Padat & gemini-embedding-001 & 3.072 & Representasi semantik dari model \textit{embedding} \textcite{googleGeminiEmbedding0012025}. \\
         & text-embedding-3-large & 3.072 & Representasi semantik dari model \textit{embedding} \textcite{openaiTextEmbedding3Large2024}. \\
        \SetCell[r=2]{l} Vektor Jarang & splade & Variabel & Representasi leksikal dari SPLADE. \\
         & bm25 & Variabel & Skor relevansi berbasis frekuensi kata (BM25). \\
        Multi-Vektor & jina-colbert-v2 & 128 & Representasi \textit{late interaction} dari ColBERT. \\
        \SetCell[r=5]{l} Payload & \textit{page\_content} & String & Konten teks lengkap dari dokumen. \\
         & \textit{metadata.plant} & String & Nama tanaman. Diindeks secara \textit{full-text}. \\
         & \textit{metadata.disease} & String & Nama penyakit. Diindeks secara \textit{full-text}. \\
         & \textit{metadata.type} & String & Jenis dokumen (budidaya/penyakit). \\
         & \textit{metadata.source} & String & Sumber dokumen. \\
        \end{longtblr}
        }

        Serupa dengan alat pengambilan citra-teks, diterapkan pula indeks teks penuh pada atribut \textit{metadata.plant} dan \textit{metadata.disease} untuk memungkinkan penyaringan berbasis token kata terhadap ruang pengambilan, sehingga hasil yang diperoleh lebih relevan dengan konteks tanaman atau penyakit yang sedang dianalisis.

        \item Teknik Pengambilan

        Dalam tahap pengambilan, dirancang teknik pencarian hibrida yang menggabungkan vektor padat dan vektor jarang, dengan hasil akhir yang disatukan menggunakan algoritma \textit{Reciprocal Rank Fusion} (RRF). Mekanisme penilaian skor RRF dirumuskan pada Persamaan~\ref{eq:rrf}.

        \begin{equation}
            score(d \in D) = \sum_{r_d \in R(d)} \frac{1}{k + r_d}
            \label{eq:rrf}
        \end{equation}
        \daftarpersamaan{Reciprocal Rank Fusion (RRF)}
        \noindent
    dengan $d$ adalah dokumen kandidat, $D$ adalah himpunan seluruh dokumen kandidat, $R(d)$ adalah himpunan peringkat dokumen $d$ dari berbagai metode pengambilan, $r_d$ adalah posisi peringkat dokumen $d$ pada suatu metode pengambilan, dan $k$ adalah konstanta penghalusan. Penerapan rumus ini diabstraksi oleh Qdrant

        Teknik pengambilan pada alat ini dirancang bertingkat, diawali dengan pemberian akses kepada agen untuk menerapkan filter metadata, seperti nama tanaman, nama penyakit, atau jenis dokumen guna mempersempit ruang pengambilan. Selanjutnya, sistem melakukan pencarian hibrida yang menggabungkan hasil dari representasi vektor padat dan vektor jarang. Sistem dikonfigurasi agar agen dapat melakukan pengambilan awal (\textit{prefetching}) terhadap sejumlah $2k$ kandidat dokumen dari masing-masing indeks vektor. Hasil dari kedua jalur pengambilan ini kemudian disatukan menggunakan algoritma RRF yang memberikan bobot berdasarkan peringkat kemunculan dokumen di setiap daftar hasil, dengan tujuan memprioritaskan dokumen yang secara konsisten dinilai relevan oleh kedua metode. 

        \noindent
        \begin{minipage}{\linewidth}
            \vspace{0.4cm}
            \captionsetup{type=figure}
            \centering
            \includegraphics[width=0.8\textwidth]{bab3/teknik_pengambilan_hibrida.png}
            \caption{Alur kerja teknik pencarian hibrida pada alat pengambilan basis pengetahuan}
            \label{fig:knowledge_base_search_workflow}
            \vspace{0.6cm}
        \end{minipage}
    \end{enumerate}

    \item Alat Pencarian Internet

    Alat pencarian internet diimplementasikan menggunakan Tavily Search API, yaitu mesin pencari yang dioptimasi khusus untuk sistem agentik dengan tujuan memberikan hasil pencarian yang efisien, cepat, dan konsisten \parencite{tavilyaiTavilyaiTavilypython2026}. Penggunaan API ini memungkinkan sistem untuk mengabstraksi serangkaian proses kompleks dalam pengambilan informasi daring, mulai dari pencarian, \textit{web scraping}, penyaringan konten, \textit{reranking}, hingga ekstraksi informasi paling relevan secara otomatis. Konfigurasi parameter yang digunakan dirangkum pada Tabel~\ref{tab:tavily_parameters}.

    {
    \setstretch{1.0}
    \begin{longtblr}[
        caption = {Konfigurasi parameter Tavily Search API},
        label = {tab:tavily_parameters},
    ]{
        width = \linewidth,
        colspec = {X[2,l] X[1,c] X[3.5,l]},
        hlines, vlines,
        rows = {m},
        rowsep = 3pt,
        row{1} = {font=\bfseries, halign=c, valign=m},
        column{1} = {font=\bfseries},
        rowhead = 1,
    }
    Parameter & Nilai & Deskripsi \\
    \textit{max\_results} & 5 & Jumlah maksimal hasil dokumen yang dikembalikan. \\
    \textit{topic} & "general" & Kategori topik pencarian untuk cakupan sumber umum. \\
    \textit{search\_depth} & "advanced" & Kedalaman pencarian untuk hasil yang lebih komprehensif. \\
    \textit{include\_images} & \textit{True} & Mengembalikan tautan citra yang relevan. \\
    \textit{include\_image\allowbreak\_descriptions} & \textit{True} & Menyertakan deskripsi tekstual untuk citra yang ditemukan. \\
    \end{longtblr}
    }

    Parameter \textit{max\_results} diatur untuk mengembalikan lima dokumen teratas. Pembatasan ini bertujuan untuk menyeimbangkan kebutuhan akan variasi sumber informasi bagi agen dengan efisiensi penggunaan API, serta mencegah masuknya informasi yang kurang relevan (\textit{noise}). Topik pencarian diatur ke "general" untuk memastikan cakupan sumber yang luas, dan fitur pencarian citra diaktifkan untuk memperkaya konteks visual yang dapat diberikan kepada pengguna.

\end{enumerate}

\subsubsection{\textit{Construction}}
Tahapan \textit{construction} oleh \textcite{pressmanSoftwareEngineeringPractitioners2020} mencakup pembuatan kode atau implementasi beserta pengujian yang diperlukan. Proses pengujian yang dilakukan mencakup dua pendekatan. Alat pengambilan citra-teks dan alat pencarian internet diuji secara \textit{black box testing} untuk memvalidasi fungsionalitas masing-masing alat, sedangkan alat pengambilan basis data pengetahuan diuji menggunakan metrik kuantitatif standar pengambilan informasi guna menentukan konfigurasi \textit{hybrid search} dan \textit{reranker} yang optimal.

\subsubsection{\textit{Deployment}}
Menurut \textcite{pressmanSoftwareEngineeringPractitioners2020}, tahap \textit{deployment} dalam pengembangan perangkat lunak merupakan fase \textit{deployment} (penerapan) perangkat lunak yang telah selesai kepada pengguna akhir. Dalam konteks pengembangan modul pengambilan informasi, tahap ini difokuskan pada penerapan basis data vektor yang menyimpan indeks citra-teks dan basis pengetahuan ke lingkungan produksi yang stabil dan dapat diakses di internet.

Selama fase pengembangan dan pengujian awal, koleksi basis data vektor untuk alat pengambilan citra-teks maupun alat pengambilan basis pengetahuan di-\textit{deploy} secara lokal menggunakan \textit{image} Docker resmi yang disediakan oleh Qdrant. Pendekatan ini memungkinkan iterasi cepat dan \textit{debugging} tanpa ketergantungan pada koneksi internet. Akan tetapi, setelah integrasi sistem secara keseluruhan, infrastruktur penyimpanan basis data vektor perlu dipisahkan dari \textit{container} aplikasi sistem agentik. Hal ini dikarenakan tempat \textit{deployment} yang digunakan untuk sistem agentik tidak dirancang untuk menangani persistensi data skala besar dan beban kerja basis data vektor secara efisien.

Oleh karena itu, layanan Qdrant Cloud dipilih sebagai solusi \textit{deployment} basis data vektor. Pemilihan ini didasarkan pada kompatibilitas penuh dengan implementasi lokal yang telah dibangun menggunakan Qdrant, serta arsitektur \textit{serverless} yang membebaskan pengembangan dari kompleksitas konfigurasi dan pemeliharaan \textit{server}. Proses \textit{deployment} dilakukan dengan membuat sebuah \textit{cluster} di Qdrant Cloud. \textit{Cluster} dalam konteks ini merujuk pada instansi basis data terkelola yang dialokasikan sumber daya komputasi dan penyimpanannya secara terisolasi di infrastruktur \textit{cloud} (awan). \textit{Cluster} ini dikonfigurasi menggunakan paket \textit{Free Tier} yang disediakan oleh Qdrant Cloud, yang dinilai memadai untuk menampung volume data dan beban \textit{query} dalam lingkup penelitian ini. Konfigurasi \textit{cluster} yang digunakan diperlihatkan pada Gambar \ref{fig:qdrant_cloud_config}.

\begin{figure}[H]
    \centering
    \includegraphics[width=0.8\textwidth]{bab3/qdrant_cloud_config.png}
    \caption{Konfigurasi \textit{cluster} di Qdrant Cloud}
    \label{fig:qdrant_cloud_config}
\end{figure}

Adapun komponen perangkat lunak sisanya dari modul ini, yakni kode program dan model-model yang menangani logika pengambilan informasi dan integrasi dengan layanan eksternal, dikemas dan di-\textit{deploy} menggunakan \textit{image} Docker yang sama dengan komponen sistem agentik lainnya. Serupa dengan modul deteksi objek, model-model yang digunakan dalam modul ini juga dikonversi ke format ONNX dan dijalankan menggunakan ONNX Runtime dengan konfigurasi CPU.

\subsection{Antarmuka Sistem Agentik}
Antarmuka sistem dikembangkan sebagai media interaksi utama antara pengguna dan sistem agentik yang telah dibangun. Pengembangan ini bertujuan untuk memfasilitasi akses terhadap fitur identifikasi penyakit tanaman secara intuitif melalui platform berbasis web. Sejalan dengan modul pengambilan informasi, pengembangan antarmuka sistem mengikuti tahapan SDLC dengan model \textit{waterfall} berdasarkan perumusan \textcite{pressmanSoftwareEngineeringPractitioners2020}. Rincian pelaksanaan setiap tahapan dijabarkan sebagai berikut.

\subsubsection{\textit{Communication}}
Tahap \textit{communication} bertujuan memahami kebutuhan teknis yang diperlukan sebelum pengembangan antarmuka dimulai. Antarmuka sistem mencakup dua komponen utama, yaitu aplikasi web sebagai \textit{frontend} dan infrastruktur \textit{server} sebagai \textit{backend}. Infrastruktur \textit{server} mencakup seluruh komponen yang beroperasi dan berinteraksi dengan sistem agentik, termasuk model-model inferensi pada setiap modul. Seluruh komponen tersebut perlu diintegrasikan dalam satu lingkungan terpadu dan terisolasi guna menyederhanakan proses penerapan serta memastikan konsistensi operasional antar-lingkungan pengembangan dan produksi. Komponen \textit{frontend} memerlukan kerangka kerja pengembangan web yang mendukung \textit{rendering} sisi klien yang interaktif dan mekanisme \textit{streaming} untuk menampilkan respons agen secara waktu nyata. Komponen \textit{backend} memerlukan dukungan infrastruktur \textit{server} yang mampu menyediakan antarmuka pemrograman aplikasi atau \textit{Application Programming Interface} (API) untuk memfasilitasi komunikasi asinkron, manajemen \textit{thread}, serta persistensi status percakapan bagi aplikasi sistem agentik yang dikembangkan dengan LangGraph. 

\subsubsection{\textit{Planning}}
Tahap \textit{planning} bertujuan mendefinisikan sumber daya dan tumpukan teknologi yang akan digunakan sebelum pengembangan antarmuka dimulai. Berdasarkan kebutuhan yang diidentifikasi pada tahap \textit{communication}, ditetapkan kerangka kerja utama untuk masing-masing komponen antarmuka. Selain itu, ditetapkan pula strategi pengemasan komponen \textit{backend} untuk memenuhi kebutuhan lingkungan yang terpadu dan terisolasi.

Untuk komponen \textit{frontend}, dipilih Next.js \parencite{vercel_nextjs} sebagai kerangka kerja pengembangan web. Pemilihan ini didasarkan pada dukungannya terhadap \textit{rendering} sisi klien yang interaktif, ekosistem React yang matang, serta kompatibilitasnya dengan Agent Chat UI dari \textcite{langchainLangchainaiAgentchatui2026} yang diadopsi sebagai basis antarmuka percakapan. Penggunaan Agent Chat UI memungkinkan pengembangan dipercepat karena komponen inti antarmuka seperti \textit{streaming} respons dan visualisasi pemanggilan alat sudah tersedia. Untuk komponen \textit{backend}, dipilih Aegra \parencite{ibbybuilds_aegra} sebagai infrastruktur \textit{server} untuk menjalankan sistem agentik berbasis LangGraph. Aegra merupakan \textit{server} Agent Protocol \textit{open source} yang dapat dioperasikan secara mandiri di infrastruktur sendiri, sekaligus berfungsi sebagai pengganti langsung LangSmith Deployment dengan antarmuka SDK dan API yang identik. Pemilihan Aegra didasarkan pada beberapa pertimbangan, yaitu infrastrukturnya yang dirancang khusus untuk menjalankan agen LangGraph secara berstatus (\textit{stateful}), termasuk manajemen \textit{thread} dan persistensi status percakapan. Untuk lapisan persistensi data ini, digunakan layanan basis data PostgreSQL dari Supabase \parencite{supabase_github} yang konfigurasi dan koneksinya dikelola serta diabstraksi secara otomatis oleh Aegra. Pertimbangan lainnya meliputi kemampuan Aegra untuk di-\textit{deploy} secara mandiri tanpa ketergantungan pada platform komersial berbayar, serta kompatibilitasnya langsung dengan Agent Chat UI yang digunakan pada komponen \textit{frontend}. 

Untuk memenuhi kebutuhan lingkungan yang terpadu dan terisolasi, seluruh komponen \textit{backend} beserta model-model inferensi dari setiap modul dikemas dalam satu \textit{image} Docker. Pendekatan ini memastikan konsistensi dependensi dan konfigurasi antar-lingkungan pengembangan dan produksi. \textit{Image} Docker ini kemudian di-\textit{deploy} pada layanan \textit{cloud} yang mendukung aplikasi berbasis Docker.

\subsubsection{\textit{Modeling}}
Tahap \textit{modeling} mencakup perancangan arsitektur antarmuka, pemodelan fungsionalitas sistem melalui diagram, dan perancangan prototipe tampilan. Rincian setiap aspek pemodelan diuraikan sebagai berikut.

\begin{enumerate}[label=\alph*., listparindent=\parindent, labelindent=0pt]
    \item Arsitektur Antarmuka

    Sistem antarmuka dibangun dengan memisahkan komponen sisi klien (\textit{frontend}) dan sisi \textit{server} (\textit{backend}) guna memastikan skalabilitas dan kemudahan pemeliharaan.     Visualisasi arsitektur sistem secara garis besar disajikan pada Gambar \ref{fig:high_level_architecture}. Arsitektur sistem secara keseluruhan terdiri dari dua komponen utama:
    \begin{enumerate}[label=\arabic*., listparindent=\parindent, labelindent=0pt]
        \item Aplikasi Web (\textit{Frontend}): Komponen ini merupakan aplikasi web berbasis Next.js yang mengadopsi Agent Chat UI dari \textcite{langchainLangchainaiAgentchatui2026}. Aplikasi ini bertanggung jawab untuk merender antarmuka percakapan, mengelola unggahan berkas, serta menampilkan visualisasi pemanggilan alat (\textit{tool calls}) secara waktu nyata.
        \item Aegra (\textit{Backend}): Komponen ini berfungsi sebagai infrastruktur \textit{server} Agent Protocol yang menjalankan logika agen LangGraph secara berstatus (\textit{stateful}). Server menyediakan API yang dioptimalkan untuk komunikasi asinkron, manajemen \textit{thread}, dan persistensi status percakapan melalui PostgreSQL. Komponen sistem agentik lainnya dikemas dalam kontainer Docker yang terintegrasi dalam satu lingkungan.
    \end{enumerate}

    \noindent
    \begin{minipage}{\linewidth}
        \vspace{0.4cm}
        \captionsetup{type=figure}
        \centering
        \includegraphics[width=0.7\textwidth]{bab3/arsitektur_antarmuka_sistem.png}
        \caption{Arsitektur antarmuka sistem agentik}
        \label{fig:high_level_architecture}
        \vspace{0.6cm}
    \end{minipage}

    \item \textit{Use Case Diagram}

    Fungsionalitas antarmuka sistem agentik dipetakan melalui \textit{use case diagram} untuk memberikan gambaran mengenai interaksi aktor dengan sistem. Aktor utama dalam sistem ini adalah Pengguna yang berinteraksi langsung dengan aplikasi web. Gambar \ref{fig:use_case_diagram} menyajikan pemetaan fungsionalitas utama yang tersedia pada antarmuka.    Penjelasan dari fungsionalitas tersebut adalah sebagai berikut:
    \begin{enumerate}[label=\arabic*., listparindent=\parindent, labelindent=0pt]
        \item Interaksi \textit{Multimodal}: Pengguna dapat mengirimkan instruksi berupa teks (seperti pertanyaan tentang gejala) serta mengunggah citra (daun, buah, atau bagian tanaman lainnya) sebagai masukan utama proses identifikasi.
        \item Visualisasi Hasil: Sistem menampilkan respons AI secara \textit{streaming} beserta keluaran setiap alat ke pengguna.
        \item Transparansi Agen: Pengguna dapat memantau langkah-langkah internal agen (penggunaan alat deteksi, pencarian web, dll.) secara transparan untuk memverifikasi proses pengambilan keputusan.
        \item Manajemen Konteks: Pengguna dapat menyimpan, mengganti nama, atau menghapus sesi percakapan untuk mengelola riwayat identifikasi yang telah dilakukan sebelumnya.
    \end{enumerate}

    \noindent
    \begin{minipage}{\linewidth}
        \vspace{0.4cm}
        \captionsetup{type=figure}
        \centering
        \includegraphics[width=0.5\textwidth]{bab3/usecasediagram.png}
        \caption{\textit{Use case diagram} antarmuka sistem agentik}
        \label{fig:use_case_diagram}
        \vspace{0.6cm}
    \end{minipage}

    \item \textit{Sequence Diagram}

    Alur komunikasi dan koordinasi antar komponen dalam sistem agentik dipetakan melalui \textit{sequence diagram}. Diagram ini memberikan gambaran mengenai urutan pertukaran pesan dan data antar aktor serta komponen sistem dalam satu skenario tugas identifikasi penyakit tanaman. Gambar \ref{fig:sequence_diagram} memvisualisasikan urutan interaksi sistem dari masukan pengguna hingga pemberian identifikasi akhir.

    \noindent
    \begin{minipage}{\linewidth}
        \vspace{0.4cm}
        \captionsetup{type=figure}
        \centering
        \includegraphics[width=1\textwidth]{bab3/sequencec_diagram_identifikasi.png}
        \caption{\textit{Sequence diagram} interaksi sistem agentik dalam identifikasi penyakit tanaman}
        \label{fig:sequence_diagram}
        \vspace{0.6cm}
    \end{minipage}

    Secara lebih rinci, arsitektur sistem terdiri dari tiga lapisan utama yang bekerja sama untuk memfasilitasi interaksi pengguna dengan agen cerdas, yaitu:
    \begin{enumerate}[nosep]
        \item \textit{Frontend}: Bertindak sebagai antarmuka pengguna yang menangani tampilan visual, penyimpanan kunci API secara lokal, manajemen \textit{state} sementara dengan React, serta pengiriman permintaan ke \textit{server}.
        \item Aegra: Berperan sebagai orkestrator yang mengelola siklus hidup \textit{thread} (kontainer percakapan) dan \textit{run} (eksekusi), menyediakan titik akhir API, serta menangani antrian eksekusi untuk \textit{background worker}.
        \item Sistem Agentik: Menjalankan logika bisnis utama dan mengeksekusi graf, termasuk pemanggilan alat eksternal seperti basis data dan model AI untuk menghasilkan respons.
    \end{enumerate}

    Interaksi antara ketiga komponen tersebut dibagi menjadi tiga skenario utama, yaitu pengiriman pesan baru, akses \textit{thread} lama, dan pengiriman pesan lanjutan.

    Skenario pertama adalah pengiriman pesan baru, baik pada \textit{thread} baru maupun \textit{thread} yang sudah ada. Sebagaimana diperlihatkan pada Gambar \ref{fig:seq_submit}, proses dimulai ketika pengguna mengirimkan pesan yang kemudian diteruskan oleh \textit{frontend} ke API untuk membuat \textit{run} baru. \textit{Run} tersebut dimasukkan ke dalam antrian untuk diproses oleh \textit{background worker}. Sistem agentik kemudian mengeksekusi graf secara lengkap, mulai dari pengambilan citra, inferensi model, hingga pencarian vektor, dan mengalirkan hasilnya (\textit{streaming}) kembali ke antarmuka pengguna secara waktu nyata.

    \noindent
    \begin{minipage}{\linewidth}
        \vspace{0.4cm}
        \captionsetup{type=figure}
        \centering
        \includegraphics[width=1\textwidth]{bab3/sequencediagramsubmit.png}
        \caption{\textit{sequence diagram} pengiriman pesan baru}
        \label{fig:seq_submit}
        \vspace{0.6cm}
    \end{minipage}

    Skenario kedua menjelaskan mekanisme akses ke \textit{thread} lama melalui riwayat percakapan. Pada skenario ini, yang divisualisasikan pada Gambar \ref{fig:seq_access_thread}, tidak terjadi panggilan API ke Aegra untuk eksekusi graf. \textit{Frontend} hanya melakukan resolusi konfigurasi (seperti ID asisten dan URL API) dan menyinkronkan \textit{state} dengan parameter URL. Riwayat percakapan ditampilkan menggunakan data yang telah tersedia tanpa memicu proses komputasi ulang di sisi agen.

    \noindent
    \begin{minipage}{\linewidth}
        \vspace{0.4cm}
        \captionsetup{type=figure}
        \centering
        \includegraphics[width=1\textwidth]{bab3/sequencediagramaksesthreadlama.png}
        \caption{\textit{Sequence diagram} akses \textit{thread} sebelumnya}
        \label{fig:seq_access_thread}
        \vspace{0.6cm}
    \end{minipage}

    Skenario ketiga menunjukkan alur ketika pengguna melanjutkan percakapan pada \textit{thread} yang sama (\textit{resubmit}), sebagaimana ditunjukkan pada Gambar \ref{fig:seq_resubmit}. Dalam kondisi ini, \textit{server} membuat \textit{run} baru yang dikonfigurasi untuk menggunakan status (\textit{state}) dari \textit{thread} sebelumnya sebagai konteks, sehingga agen dapat memproses pesan baru dengan tetap mempertahankan memori dari interaksi sebelumnya.

    \noindent
    \begin{minipage}{\linewidth}
        \vspace{0.4cm}
        \captionsetup{type=figure}
        \centering
        \includegraphics[width=1\textwidth]{bab3/submitulangdithreadlama.png}
        \caption{\textit{Sequence diagram} melanjutkan percakapan pada \textit{thread} lama}
        \label{fig:seq_resubmit}
        \vspace{0.6cm}
    \end{minipage}

    \item Prototipe Antarmuka

    Antarmuka pada sistem ini diadopsi dari proyek sumber terbuka Agent Chat UI yang dikembangkan oleh \textcite{langchainLangchainaiAgentchatui2026}. Desain antarmuka dirancang dengan pendekatan minimalis untuk memfokuskan perhatian pengguna pada interaksi percakapan dan hasil identifikasi. Gambar \ref{fig:prototipe_antarmuka} memperlihatkan tampilan utama dari prototipe antarmuka. 
    Pada bagian kiri antarmuka terdapat panel sisi (\textit{sidebar}) yang menampilkan riwayat percakapan (\textit{Thread History}) pengguna. Bagian tengah antarmuka didominasi oleh area interaksi utama yang memuat logo sistem, kolom input pesan, dan tombol aksi. Kolom input dilengkapi dengan fitur untuk mengunggah berkas citra, serta tombol geser (\textit{toggle}) untuk menyembunyikan atau menampilkan detail pemanggilan alat (\textit{tool calls}) oleh agen.

    \noindent
    \begin{minipage}{\linewidth}
        \vspace{0.4cm}
        \captionsetup{type=figure}
        \centering
        \includegraphics[width=1\textwidth]{bab3/prototipeantarmuka.png}
        \caption{Tampilan prototipe antarmuka sistem agentik}
        \label{fig:prototipe_antarmuka}
        \vspace{0.6cm}
    \end{minipage}
\end{enumerate}

\subsubsection{\textit{Construction}}
Tahap \textit{construction} mencakup implementasi komponen antarmuka berdasarkan perancangan yang telah dilakukan dan pengujian fungsionalitasnya. Pengujian antarmuka dilakukan menggunakan metode \textit{black box} untuk memastikan seluruh fungsionalitas yang dirancang pada Gambar \ref{fig:use_case_diagram} dapat beroperasi sesuai dengan ekspektasi pengguna. Rincian skenario dan prosedur pengujian diuraikan pada bagian Pengujian.

\subsubsection{\textit{Deployment}}
Menurut \textcite{pressmanSoftwareEngineeringPractitioners2020}, tahap \textit{deployment} merupakan fase penerapan perangkat lunak yang telah selesai kepada pengguna akhir. Komponen \textit{frontend} berupa aplikasi web Next.js di-\textit{hosting} pada platform Vercel \parencite{vercel_hosting}. Komponen \textit{backend} berupa Aegra dikemas dalam kontainer dan diterapkan pada layanan \textit{cloud container} Google Cloud Run. Seluruh komponen lokal lainnya, dikemas dalam satu \textit{image} Docker yang sama dan di-\textit{deploy} bersama dengan layanan sistem agentik. 

\section{Pengujian}
Tahap pengujian dilakukan untuk memverifikasi bahwa komponen-komponen sistem berfungsi sesuai spesifikasi dan terintegrasi dengan benar ke dalam kerangka sistem agentik.

\subsection{Modul Deteksi Objek}
Pengujian modul deteksi objek tanaman mencakup dua aspek, yaitu pengujian integrasi fungsional dan pengukuran latensi inferensi. Pengujian integrasi dilakukan menggunakan pendekatan \textit{black box testing} sesuai dengan \textcite{bhasinBlackBoxTesting2014}. Pengujian ini mencakup kedua alat dalam modul, yaitu alat deteksi \textit{closed-set} daun dan alat \textit{open-vocabulary object detection}, untuk memastikan masing-masing alat dapat dipanggil oleh agen, memproses citra masukan, dan mengembalikan keluaran deteksi sesuai skema yang telah didefinisikan. Rincian skenario pengujian \textit{black box} untuk modul ini disajikan pada Tabel \ref{tab:black_box_testing_objek}.

{
\setstretch{1.0}
\begin{longtblr}[
    caption = {Pengujian \textit{black box} integrasi modul deteksi objek},
    entry = {Pengujian \textit{black box} Integrasi Modul Deteksi Objek},
    label = {tab:black_box_testing_objek},
]{
    width = \linewidth,
    colspec = {X[1.5,l] X[2.3,l] X[2.3,l]},
    hlines, vlines,
    rows = {m},
    rowhead = 1,
    row{1} = {font=\bfseries, halign=c, valign=m},
    rowsep = 3pt,
}
    Fungsionalitas & Skenario Pengujian & Hasil yang Diharapkan \\
    Pemanggilan Manual Alat Deteksi \textit{Closed-Set} Daun & Menjalankan fungsi alat secara terisolasi dengan input citra tanaman. & Model YOLOv11 berhasil melakukan inferensi dan mengembalikan koordinat \textit{bounding box} daun. \\
    Pemanggilan Manual Alat OVD & Menjalankan fungsi alat secara terisolasi dengan input citra dan \textit{text prompt} target. & Model OWLv2 berhasil melakukan inferensi dan mengembalikan koordinat objek sesuai \textit{prompt}. \\
    Integrasi Agen untuk Alat Deteksi \textit{Closed-Set} Daun & Mengirimkan instruksi eksplisit kepada agen untuk memanggil alat deteksi \textit{closed-set} daun pada citra daun. & Agen mengenali instruksi spesifik, memanggil alat deteksi \textit{closed-set} daun, dan melaporkan keluaran alat. \\
    Integrasi Agen untuk Alat OVD & Mengirimkan instruksi eksplisit kepada agen untuk memanggil alat OVD pada bagian tanaman tertentu dari citra. & Agen mengenali instruksi spesifik, memanggil alat OVD dengan \textit{prompt} yang sesuai, dan melaporkan keluaran alat. \\
\end{longtblr}
}

Selain pengujian integrasi, dilakukan juga \textit{benchmarking} latensi inferensi yang difokuskan pada varian model YOLOv11 guna menentukan model terbaik yang akan digunakan oleh alat deteksi \textit{closed-set} daun. Pengujian ini bertujuan untuk mengukur efisiensi model saat dijalankan pada perangkat dengan sumber daya terbatas. Prosedur pengujian terdiri dari dua tahap, yaitu fase pemanasan (\textit{warm-up}) dan fase pengujian utama. Fase pemanasan dilakukan sebanyak 30 iterasi untuk menstabilkan performa sistem dan memastikan model telah dimuat sepenuhnya ke dalam memori. Selanjutnya, fase pengujian utama dilakukan sebanyak 300 iterasi menggunakan citra yang dipilih secara acak dari dataset. 

Pada setiap iterasi, waktu eksekusi dicatat untuk tiga komponen utama, yaitu prapemrosesan, inferensi model, dan pascapemrosesan. Prapemrosesan meliputi proses \textit{resizing} citra ke resolusi $416 \times 416$ piksel, normalisasi, dan penyusunan \textit{batch}. Tahap inferensi model dijalankan untuk menghasilkan prediksi deteksi. Tahap pascapemrosesan meliputi penskalaan koordinat dan penerapan NMS untuk menyaring hasil deteksi. Latensi total dihitung sebagai akumulasi dari ketiga tahapan tersebut. Alur kerja prosedur \textit{benchmarking} ini diilustrasikan pada Gambar \ref{fig:benchmark_flowchart}.

\begin{figure}[H]
    \centering
    \includegraphics[width=1\textwidth]{bab3/flowchart_benchmark_yolo.png}
    \caption{Diagram alur prosedur \textit{benchmarking} latensi model}
    \label{fig:benchmark_flowchart}
\end{figure}

Data latensi yang dikumpulkan kemudian dianalisis menggunakan metrik persentil ke-50 (p50) atau median dan persentil ke-95 (p95). Misalkan $T = \{t_1, t_2, \dots, t_n\}$ adalah himpunan waktu latensi total yang telah diurutkan secara menaik. Hasil \textit{benchmarking} dari ketiga varian YOLOv11 selanjutnya digunakan dalam tahap evaluasi untuk menentukan model yang paling optimal. Definisi dari kedua metrik latensi sendiri dijabarkan sebagai berikut:

\begin{enumerate}[nosep]
    \item p50 (Median): Merepresentasikan nilai tengah dari distribusi latensi yang menggambarkan kecepatan respons tipikal dalam penggunaan normal. Secara matematis, p50 dapat dihitung menggunakan Persamaan \ref{eq:p50}:
    \begin{equation}
        p50 = \begin{cases} t_{\frac{n+1}{2}} & \text{jika } n \text{ ganjil} \\ \frac{1}{2} (t_{\frac{n}{2}} + t_{\frac{n}{2}+1}) & \text{jika } n \text{ genap} \end{cases}
        \label{eq:p50}
    \end{equation}
    Dalam persamaan tersebut, $n$ adalah jumlah total sampel data, sedangkan $t_i$ merepresentasikan nilai latensi pada urutan ke-$i$ setelah data diurutkan secara menaik.
    \daftarpersamaan{Latensi p50 (Median)}

    \item p95 (Persentil ke-95): Merepresentasikan ambang batas waktu di mana 95\% dari total sampel inferensi lebih cepat dari nilai tersebut. Metrik ini mengukur performa dalam skenario terburuk (\textit{tail latency}), memastikan sistem tetap responsif bahkan saat memproses citra yang kompleks. Nilai p95 dihitung sesuai Persamaan \ref{eq:p95}:
    \begin{equation}
        p95 = t_{\lceil 0.95 \cdot n \rceil}
        \label{eq:p95}
    \end{equation}
    Dalam persamaan tersebut, notasi $\lceil \cdot \rceil$ melambangkan fungsi \textit{ceiling} untuk pembulatan ke atas ke bilangan bulat terdekat, dan $n$ adalah jumlah total sampel data.
    \daftarpersamaan{Latensi p95 (Persentil ke-95)}
\end{enumerate}

\subsection{Modul Pengambilan Informasi}
Pengujian pada modul pengambilan informasi dibedakan berdasarkan kebutuhan masing-masing alat. Alat pengambilan citra-teks dan alat pencarian internet tidak memerlukan \textit{benchmarking} kuantitatif sehingga diuji menggunakan pendekatan \textit{black box testing}. Pengujian ini dilakukan untuk memvalidasi fungsionalitas alat secara mandiri maupun integrasinya dengan agen. Rincian skenario pengujian \textit{black box} untuk modul ini disajikan pada Tabel \ref{tab:black_box_testing_retrieval}.

{
\setstretch{1.0}
\begin{longtblr}[
    caption = {Pengujian \textit{black box} fungsionalitas dan integrasi modul pengambilan informasi},
    entry = {Pengujian \textit{black box} Fungsionalitas dan Integrasi Modul Pengambilan Informasi},
    label = {tab:black_box_testing_retrieval},
]{
    width = \linewidth,
    colspec = {X[1.5,l] X[2.3,l] X[2.3,l]},
    hlines, vlines,
    rows = {m},
    rowhead = 1,
    row{1} = {font=\bfseries, halign=c, valign=m},
    rowsep = 3pt,
}
    Fungsionalitas & Skenario Pengujian & Hasil yang Diharapkan \\
    Teknik \textit{Text-to-Text} & Menjalankan alat pengambilan citra-teks dengan masukan deskripsi gejala teks. & Alat mengembalikan metadata penyakit yang relevan berdasarkan kemiripan semantik teks. \\
    Teknik \textit{Text-to-Image} & Menjalankan alat pengambilan citra-teks dengan masukan deskripsi teks. & Alat mengembalikan citra referensi visual yang sesuai dengan deskripsi teks. \\
    Teknik \textit{Image-to-Image} & Menjalankan alat pengambilan citra-teks dengan masukan citra daun bergejala. & Alat mengembalikan citra penyakit dari basis data yang memiliki kemiripan visual paling tinggi. \\
    Teknik \textit{Image-to-Text} & Menjalankan alat pengambilan citra-teks dengan masukan citra. & Alat mengembalikan label atau deskripsi penyakit yang relevan dengan fitur visual citra. \\
    Filter Pengambilan Citra-Teks & Menjalankan pencarian \textit{multimodal} dengan filter metadata spesies tanaman tertentu. & Hasil pencarian hanya mencakup entri penyakit yang berasosiasi dengan spesies tanaman yang difilter. \\
    Pemanggilan Manual Alat Pengambilan Citra-Teks Tanpa Deteksi & Menjalankan fungsi alat secara terisolasi dengan input citra tanpa bantuan deteksi objek. & Alat berhasil melakukan pengambilan citra-teks pada citra penuh dan mengembalikan kandidat penyakit yang relevan. \\
    Pemanggilan Manual Alat Pengambilan Citra-Teks Dengan Deteksi & Menjalankan fungsi alat secara terisolasi dengan input citra dan hasil deteksi objek sebagai panduan area. & Alat berhasil melakukan pengambilan citra-teks pada setiap regio terdeteksi dan mengembalikan prediksi penyakit spesifik untuk masing-masing area. \\
    Pemanggilan Manual Alat Pengambilan Basis Pengetahuan Tanpa Filter & Menjalankan fungsi alat secara terisolasi dengan \textit{query} tentang penyakit tanaman tanpa parameter filter. & Alat berhasil melakukan pencarian hibrida pada basis data vektor dan mengembalikan dokumen relevan dari seluruh korpus. \\
    Pemanggilan Manual Alat Pengambilan Basis Pengetahuan Dengan Filter & Menjalankan fungsi alat secara terisolasi dengan \textit{query} dan parameter filter metadata. & Hasil dokumen yang dikembalikan relevan dengan \textit{query} dan terbatas pada entri yang sesuai dengan kriteria filter metadata. \\
    Pemanggilan Manual Alat Pencarian Internet & Menjalankan fungsi alat secara terisolasi dengan \textit{query} informasi umum pertanian. & Tavily API berhasil mengembalikan daftar hasil pencarian yang relevan beserta ringkasannya. \\
    Integrasi Agen untuk Alat Pengambilan Citra-Teks & Mengirimkan citra tanaman bergejala kepada agen untuk identifikasi. & Agen memanggil alat pengambilan citra-teks secara otomatis dan menggunakan hasilnya untuk memberikan hipotesis identifikasi. \\
    Integrasi Agen untuk Alat Pengambilan Basis Pengetahuan & Mengajukan pertanyaan spesifik mengenai manajemen penyakit yang tercakup dalam korpus. & Agen memanggil alat pengambilan basis pengetahuan dan memberikan panduan yang didukung oleh kutipan dari dokumen korpus. \\
        Integrasi Agen untuk Alat Pencarian Internet & Mengajukan pertanyaan mengenai informasi eksternal yang dinamis. & Agen mengenali kebutuhan informasi eksternal, memanggil alat pencarian internet, dan menyajikan jawaban berdasarkan hasil pencarian. \\
\end{longtblr}
}

Sementara itu, alat pengambilan basis data pengetahuan memerlukan pengujian kuantitatif kinerja secara \textit{offline} menggunakan pustaka ranx oleh \textcite{bassaniRanxBlazingFastPython2022}. Pengujian ini bertujuan untuk mengukur efektivitas berbagai teknik pengambilan dan konfigurasi model yang telah diimplementasikan dalam menemukan dokumen relevan dari korpus basis data pengetahuan.

Pengujian kuantitatif difokuskan pada konfigurasi pencarian hibrida, yaitu metode yang menggabungkan vektor padat untuk menangkap konteks semantik dan vektor jarang untuk pencocokan kata kunci yang presisi. Pada komponen vektor jarang, dua pendekatan yang diuji adalah algoritma BM25 \parencite{robertsonProbabilisticRelevanceFramework2009} yang berbasis frekuensi kata dan tersedia di Qdrant, serta model SPLADE \parencite{formalSPLADESparseLexical2021} yang berbasis pembelajaran mesin dan dijalankan melalui FastEmbed dari \textcite{qdrant_fastembed_2023}. Kombinasi keduanya dengan dua model \textit{embedding} padat menghasilkan empat variasi konfigurasi yang diuji, yaitu text-embedding-3-large dengan BM25, text-embedding-3-large dengan SPLADE, gemini-embedding-001 dengan BM25, dan gemini-embedding-001 dengan SPLADE.

Konfigurasi pencarian hibrida terbaik selanjutnya diuji lebih lanjut dengan menambahkan komponen \textit{reranker}, untuk membandingkan dua arsitektur yang berbeda, yakni model berbasis \textit{cross-encoder} dan model dengan pendekatan \textit{late interaction}. Sebagai representasi \textit{cross-encoder}, dipilih model rerank-2.5 dari Voyage AI yang memproses \textit{query} dan dokumen secara bersamaan untuk akurasi tinggi. Untuk pendekatan \textit{late interaction}, digunakan model jina-colbert-v2 dari \textcite{jhaJinaColBERTv2GeneralPurposeMultilingual2024} yang diimplementasikan melalui FastEmbed. Keseluruhan alur pengujian komponen pengambilan ini divisualisasikan pada Gambar~\ref{fig:alur_pengujian_retrieval}.

\begin{figure}[H]
    \centering
    \includegraphics[width=\textwidth]{bab3/flowchart_pengujian_rag.png}
    \caption{Diagram alur pengujian komponen alat pencarian basis data pengetahuan}
    \label{fig:alur_pengujian_retrieval}
\end{figure}

Persiapan pengujian dimulai dengan penyusunan Qrels. Qrels merupakan standar kebenaran (\textit{ground truth}) yang memetakan setiap \textit{query} dengan dokumen-dokumen yang dianggap relevan. Dalam penelitian ini, struktur Qrels direpresentasikan dalam format \textit{dictionary} Python sesuai dengan spesifikasi ranx, di mana setiap ID \textit{query} memiliki daftar ID dokumen relevan beserta skor relevansinya. Skor relevansi yang digunakan bersifat biner, di mana nilai 1 diberikan untuk dokumen yang relevan dan secara implisit nilai 0 untuk dokumen lainnya. Hal ini didasarkan pada fakta bahwa dataset QA sintetis dibangun dengan memasangkan \textit{query} secara spesifik pada potongan dokumen sumber tertentu, sehingga relevansi dapat ditentukan secara deterministik. 

Proses pengujian dilakukan dengan membandingkan hasil pengambilan sistem (\textit{Run}) terhadap Qrels menggunakan serangkaian metrik standar dalam pengambilan informasi. Mengingat penggunaan model \textit{embedding} dan LLM yang mutakhir (\textit{State-of-the-Art}), evaluasi difokuskan pada peringkat atas (Top-K) yang ketat, yaitu dengan nilai $k=5$. Pemilihan $k$ yang kecil ini bertujuan untuk memastikan sistem benar-benar menyajikan informasi yang paling relevan pada urutan teratas, sehingga meminimalkan beban kognitif pengguna dalam memilah informasi. Berikut adalah metrik-metrik yang digunakan dalam evaluasi:

\begin{enumerate}[noitemsep]
    \item \textit{Mean Reciprocal Rank} (MRR) \\
    MRR mengukur rata-rata kebalikan dari peringkat dokumen relevan pertama yang ditemukan. Metrik ini sangat krusial untuk jenis pertanyaan faktual di mana pengguna hanya membutuhkan satu jawaban yang benar secepat mungkin. Nilai MRR dihitung dengan Persamaan \ref{eq:mrr}:
    \begin{equation}
        MRR = \frac{1}{|Q|} \sum_{i=1}^{|Q|} \frac{1}{rank_i}
        \label{eq:mrr}
    \end{equation}
    \daftarpersamaan{Mean Reciprocal Rank (MRR)}
    \noindent
    Dalam persamaan tersebut, $|Q|$ adalah jumlah total \textit{query} yang dievaluasi dan $rank_i$ adalah posisi peringkat dokumen relevan pertama untuk \textit{query} ke-$i$.

    \item \textit{Normalized Discounted Cumulative Gain} (NDCG@k) \\
    NDCG@k mengevaluasi kualitas pemeringkatan dengan mempertimbangkan posisi dokumen relevan dalam daftar hasil. Dokumen relevan yang muncul di posisi lebih tinggi akan memberikan kontribusi skor yang lebih besar. Metrik ini dihitung dengan membagi \textit{Discounted Cumulative Gain} (DCG) dengan \textit{Ideal DCG} (IDCG).

    Nilai DCG pada peringkat ke-$k$ didefinisikan pada Persamaan \ref{eq:dcg} sebagai akumulasi skor relevansi yang didiskon secara logaritmik:
    \begin{equation}
        DCG@k = \sum_{i=1}^{k} \frac{rel_i}{\log_2(i+1)}
        \label{eq:dcg}
    \end{equation}
    \daftarpersamaan{Discounted Cumulative Gain (DCG@k)}
    \noindent
    Dalam persamaan tersebut, $rel_i$ merepresentasikan skor relevansi dari dokumen pada peringkat ke-$i$, sedangkan $k$ adalah batas peringkat pemotongan (\textit{cutoff}) evaluasi.

    Sedangkan IDCG adalah nilai DCG maksimal yang dapat dicapai jika dokumen-dokumen relevan diurutkan pada posisi ideal (paling atas) sebagaimana ditunjukkan pada Persamaan \ref{eq:idcg}:
    \begin{equation}
        IDCG@k = \sum_{i=1}^{|REL_k|} \frac{rel_i}{\log_2(i+1)}
        \label{eq:idcg}
    \end{equation}
    \daftarpersamaan{Ideal Discounted Cumulative Gain (IDCG@k)}
    \noindent
    Dalam persamaan tersebut, $|REL_k|$ merepresentasikan jumlah dokumen relevan yang dipertimbangkan hingga peringkat ke-$k$ (nilai minimum antara $k$ dan total dokumen relevan), dan $rel_i$ adalah skor relevansi dokumen pada posisi ideal ke-$i$.

    Sehingga, nilai NDCG@k diperoleh melalui Persamaan \ref{eq:ndcg}:
    \begin{equation}
        NDCG@k = \frac{DCG@k}{IDCG@k}
        \label{eq:ndcg}
    \end{equation}
    \daftarpersamaan{Normalized Discounted Cumulative Gain (NDCG@k)}
    \noindent
    Metrik ini digunakan untuk memastikan bahwa algoritma pengambilan tidak hanya menemukan dokumen yang benar, tetapi juga mampu mengurutkannya berdasarkan tingkat kepentingannya.

    \item \textit{Recall}@k \\
    Recall@k mengukur proporsi dokumen relevan yang berhasil ditemukan dalam $k$ hasil teratas dibandingkan dengan total dokumen relevan yang tersedia untuk \textit{query} tersebut.
    \begin{equation}
        Recall@k = \frac{|D_{rel} \cap D_{k}|}{|D_{rel}|}
        \label{eq:recall_at_k}
    \end{equation}
    \daftarpersamaan{Recall@k}
    \noindent
    Dalam persamaan tersebut, $D_{rel}$ adalah himpunan seluruh dokumen relevan yang tersedia untuk \textit{query}, $D_{k}$ adalah himpunan dokumen yang berhasil diambil dalam $k$ peringkat teratas, dan notasi $|\cdot|$ menunjukkan kardinalitas atau jumlah elemen dalam himpunan. Metrik ini penting untuk mengevaluasi cakupan (\textit{coverage}) sistem dalam menyajikan berbagai perspektif informasi yang mungkin tersebar di beberapa dokumen.
\end{enumerate}

\subsection{Antarmuka Sistem Agentik}
Pengujian antarmuka dilakukan menggunakan metode \textit{black box testing} untuk memastikan bahwa seluruh fungsionalitas yang telah dirancang pada Gambar \ref{fig:use_case_diagram} dapat beroperasi sesuai dengan ekspektasi pengguna. \textit{black box testing} merupakan metode pengujian perangkat lunak yang berfokus pada spesifikasi fungsional sistem \parencite{bhasinBlackBoxTesting2014}. Pengujian fungsionalitas antarmuka dirangkum dalam Tabel \ref{tab:black_box_testing_antarmuka}.

{
\setstretch{1.0}
\begin{longtblr}[
    caption = {Pengujian \textit{black box} antarmuka sistem agentik},
    entry = {Pengujian \textit{black box} Antarmuka Sistem Agentik},
    label = {tab:black_box_testing_antarmuka},
]{
    width = \linewidth,
    colspec = {X[1.2,l] X[2.5,l] X[2.5,l]},
    hlines, vlines,
    rows = {m},
    rowhead = 1,
    row{1} = {font=\bfseries, halign=c, valign=m},
    rowsep = 3pt,
}
    Fitur & Skenario Pengujian & Hasil yang Diharapkan \\
    Mengirim Prompt & Pengguna mengetik instruksi pada kolom teks (\textit{textarea}) dan menekan tombol kirim (ikon \textit{ArrowUp}). & Sistem mengirim pesan ke \textit{server} dan menampilkan respons agen secara bertahap (\textit{streaming}). \\
    Mengunggah Citra Tanaman & Pengguna mengklik ikon \textit{Paperclip} untuk memilih citra atau menyeret berkas langsung ke area percakapan. & Citra muncul sebagai \textit{preview} di atas kolom input dan berhasil dikirim bersama pesan teks. \\
    Melihat Hasil \& Saran & Pengguna menunggu proses inferensi selesai setelah mengirimkan \textit{query} identifikasi. & Jawaban berupa identifikasi penyakit dan panduan penanganan muncul secara lengkap di area percakapan. \\
    Mengelola Riwayat Percakapan & Pengguna memilih sesi pada \textit{sidebar}, mengubah nama sesi, atau mengirim pesan baru pada sesi lama. & Riwayat dimuat, nama sesi berubah, dan agen memproses pesan baru dengan konteks sesi tersebut. \\
    Melihat Transparansi Agen & Pengguna menekan tombol kontrol alat (ikon \textit{Wrench}) untuk mengaktifkan visibilitas proses. & Detail pemanggilan alat (\textit{tool calls}) dan hasil observasi internal agen ditampilkan di antara pesan percakapan. \\
\end{longtblr}
}

\subsection{Sistem Agentik}
Setelah semua komponen berhasil terintegrasi dengan baik, sistem perlu diuji secara komprehensif. Seluruh rangkaian pengujian ini diorkestrasi menggunakan platform LangSmith dari \textcite{langchain_langsmith_home}. Platform ini difungsikan sebagai pusat manajemen pengujian yang mencakup pengelolaan \textit{prompt} dan teknik \textit{prompt engineering} untuk agen LLM, manajemen dataset pengujian, serta analisis hasil eksperimen secara terpusat. Rangkaian pengujian dibagi menjadi dua kategori utama, yaitu pengujian satu tahap untuk mengukur kapabilitas penyelesaian tugas tunggal dan pengujian multi-tahap untuk mengevaluasi konsistensi interaksi dalam percakapan berkelanjutan.

\subsubsection{Pengujian Satu Tahap}
Pengujian satu tahap difokuskan pada evaluasi kemampuan agen dalam menyelesaikan tugas identifikasi dan tanya jawab dalam satu giliran interaksi tunggal. Pengujian ini terdiri dari eksperimen VQA (\textit{Visual Question Answering}). Pengujian ini merupakan uji utama yang dilakukan untuk mengevaluasi kinerja sistem secara menyeluruh. Skenario ini dirancang untuk merepresentasikan interaksi dunia nyata di mana pengguna mengunggah citra tanaman beserta \textit{prompt} berisi pertanyaan terkait kondisinya. Dalam eksperimen ini, seluruh kapabilitas sistem akan diuji secara terintegrasi, meliputi kemampuan pemahaman konteks \textit{multimodal}, akurasi identifikasi penyakit visual, ketepatan pengambilan informasi dari basis pengetahuan, kesesuaian penggunaan alat, serta kualitas dan relevansi jawaban yang diberikan kepada pengguna.

\begin{enumerate}[label=\alph*., listparindent=\parindent, labelindent=0pt]
    \item Persiapan Dataset Pengujian Satu Tahap

    Dataset yang digunakan dalam pengujian satu tahap dikelompokkan menjadi dua jenis berdasarkan distribusi datanya, yaitu dataset \textit{In-Distribution} (ID) dan dataset \textit{Out-of-Distribution} (OOD).

    \begin{enumerate}[label=\arabic*., listparindent=\parindent, labelindent=0pt]
        \item \textit{Dataset In-Distribution} (ID)

        Dataset ID merepresentasikan data yang berada dalam cakupan distribusi pengetahuan internal sistem. Data ini diambil dari subset dataset PlantWild \parencite{weiBenchmarkingInthewildMultimodal2024} yang telah digunakan dalam penyusunan indeks pengambilan citra-teks. Sebagaimana dijelaskan pada bagian sebelumnya, dataset ini mencakup 18.542 citra dari 89 kelas kondisi kesehatan pada 13 spesies tanaman. Untuk keperluan pengujian, dilakukan pengambilan sampel acak sebanyak 3 citra dari setiap kelas, sehingga diperoleh total 267 citra uji. Penggunaan sampel ini bertujuan untuk memvalidasi kemampuan sistem dalam mengenali dan mengambil informasi yang sudah pernah dipelajari atau diindeks sebelumnya.

        \item \textit{Dataset Out-of-Distribution} (OOD)

        Dataset OOD disusun untuk menguji ketangguhan dan kemampuan generalisasi sistem dalam menangani \textit{query} yang melibatkan spesies tanaman atau jenis penyakit yang tidak terdapat dalam basis pengetahuan internal. Dataset ini dikompilasi dari berbagai dataset publik yang mencakup 8 komoditas tanaman berbeda dengan total 40 kelas penyakit dan kondisi kesehatan unik. Rincian sumber data untuk setiap komoditas disajikan pada Tabel \ref{tab:ood_datasets}.

{
\setstretch{1.0}
\begin{longtblr}[
    caption = {Rincian dataset \textit{out-of-distribution} (OOD)},
    label = {tab:ood_datasets},
]{
    width = \linewidth,
    colspec = {X[0.5,c] X[2,l] X[0.8,c] X[3,l]},
    hlines, vlines,
    rows = {m},
    rowsep = 3pt,
    row{1} = {font=\bfseries, halign=c, valign=m},
    rowhead = 1,
}
    No. & Komoditas & Kelas & Sumber Referensi \\
    1. & Cabai Merah & 5 & \textcite{mpChilliDataset2024} \\
    2. & Kacang Hitam & 5 & \textcite{talasilaBlackGramPlant2022} \\
    3. & Bayam Malabar & 5 & \textcite{sayeemIDDMSLDImageDataset2025} \\
    4. & Okra & 6 & \textcite{kumarakuruOkraDiseaseNetDataset2024} \\
    5. & Lobak & 5 & \textcite{rashidImageDatasetRadish2024} \\
    6. & Jambu Biji & 7 & \textcite{shihabImageDatasetDetection2024} \\
    7. & Gandum & 3 & \textcite{getachewWheatLeafDataset2021} \\
    8. & Bawang Bombay & 4 & \textcite{mpOnionDataset2024} \\
    \SetCell[c=2]{c} \textbf{Total} & & \textbf{40} & - \\
\end{longtblr}
}

Serupa dengan prosedur pada dataset ID, dilakukan pengambilan sampel acak sebanyak 3 citra untuk setiap kelas pada dataset OOD, menghasilkan total 120 citra uji. Citra-citra tersebut kemudian juga diunggah ke layanan penyimpanan objek Cloudflare R2 \parencite{cloudflare_r2} untuk mendapatkan URL publik. Selanjutnya, metadata untuk setiap sampel disusun dalam format JSON. Metadata ini mencakup label penyakit, status kesehatan, \textit{caption} gejala visual, serta informasi manajemen penyakit yang dikurasi dari sumber ilmiah otoritatif. Informasi manajemen penyakit ini disertakan sebagai \textit{ground truth} untuk mengevaluasi kualitas saran yang diberikan oleh agen. Contoh entri dataset evaluasi OOD disajikan pada Gambar \ref{fig:ood_json_example}.

\noindent
    \begin{minipage}{\linewidth}
        \vspace{0.4cm}
        \captionsetup{type=figure}
        \centering
        \includegraphics[width=1\textwidth]{bab3/contohentri.png}
        \caption{Contoh entri dataset evaluasi OOD dalam format JSON}
        \label{fig:ood_json_example}
        \vspace{0.6cm}
    \end{minipage}

Selanjutnya, dataset evaluasi diperkaya melalui proses menghasilkan data sintetis menggunakan LLM Gemini 2.5 Flash. Proses ini bertujuan untuk mensimulasikan variasi interaksi pengguna yang mungkin terjadi di dunia nyata. Model LLM diberikan akses terhadap citra asli beserta metadata JSON yang telah disusun sebelumnya untuk membangkitkan pasangan pertanyaan-jawaban (\textit{QA pairs}) yang kontekstual. Baik pada dataset ID maupun OOD, untuk setiap citra uji dibangkitkan tiga variasi skenario pertanyaan sebagai berikut:

\begin{enumerate}[label=\arabic*., listparindent=\parindent, labelindent=0pt]
    \item Skenario 1: Pengguna diasumsikan sebagai orang awam yang mendeskripsikan gejala visual secara samar namun menyebutkan spesies tanaman.
    \item Skenario 2: Pengguna hanya menyebutkan nama spesies tanaman tanpa memberikan deskripsi gejala spesifik.
    \item Skenario 3: Pengguna tidak memiliki pengetahuan mengenai spesies maupun gejala spesifik. Pertanyaan yang diajukan bersifat umum seperti "Kenapa tanaman saya?" atau "Apakah ini penyakit?".
\end{enumerate}

Logika yang sama juga diterapkan untuk kelas tanaman sehat. Khusus untuk Skenario 1 pada tanaman sehat, deskripsi difokuskan pada kondisi visual tanaman yang menunjukkan ciri-ciri pertumbuhan normal. Untuk memastikan kualitas data yang dihasilkan, diterapkan teknik \textit{few-shot prompting} dalam instruksi yang diberikan kepada LLM. Contoh \textit{prompt} yang digunakan untuk Skenario 1 disajikan pada Gambar \ref{fig:prompt_scenario_1}.

\noindent
    \begin{minipage}{\linewidth}
        \vspace{0.4cm}
        \captionsetup{type=figure}
        \centering
        \includegraphics[width=1\textwidth]{bab3/contohpromptpersiapandata.png}
        \caption{Contoh \textit{prompt} \textit{few-shot} untuk Skenario 1}
        \label{fig:prompt_scenario_1}
        \vspace{0.6cm}
    \end{minipage}

Proses menghasilkan dataset ini dilakukan melalui alur kerja terstruktur untuk menjamin validitas dan konsistensi format keluaran. Setiap citra diproses melalui tahapan validasi kualitas, analisis komposisi visual, penentuan logika penggunaan alat, hingga menghasilkan konten akhir. Alur instruksi yang diberikan kepada LLM dalam proses ini divisualisasikan pada Gambar \ref{fig:dataset_generation_pipeline}.

\noindent
    \begin{minipage}{\linewidth}
        \vspace{0.4cm}
        \captionsetup{type=figure}
        \centering
        \includegraphics[width=0.8\textwidth]{bab3/flowchart_generate_dataset_sistem_agentik.png}
        \caption{Alur instruksi menghasilkan dataset sintetis dengan LLM}
        \label{fig:dataset_generation_pipeline}
        \vspace{0.6cm}
    \end{minipage}

{Hasil akhir dari proses menghasilkan ini disimpan dalam format JSON yang terstruktur. Struktur ini mencakup input simulasi pengguna, luaran referensi yang diharapkan (\textit{ground truth}), serta metadata pendukung. Contoh entri dataset hasil untuk ketiga skenario disajikan pada Gambar \ref{fig:synthetic_dataset_examples}. \par}

\noindent
    \begin{minipage}{\linewidth}
        \vspace{0.4cm}
        \captionsetup{type=figure}
        \centering
        \begin{subfigure}[b]{\textwidth}
            \centering
            \includegraphics[width=\textwidth]{bab3/contohskenario1.png}
            \caption{}
        \end{subfigure}
        \hfill
        \begin{subfigure}[b]{\textwidth}
            \centering
            \includegraphics[width=\textwidth]{bab3/contohskenario2.png}
            \caption{}
        \end{subfigure}
        \hfill
        \begin{subfigure}[b]{\textwidth}
            \centering
            \includegraphics[width=\textwidth]{bab3/contohskenario3.png}
            \caption{}
        \end{subfigure}
        \caption[Contoh entri dataset evaluasi sintetis]{Contoh entri dataset evaluasi sintetis: (a) Skenario 1, (b) Skenario 2, dan (c) Skenario 3}
        \label{fig:synthetic_dataset_examples}
        \vspace{0.6cm}
    \end{minipage}
    \end{enumerate}

\item Metrik Pengukuran Pengujian Satu Tahap

Kinerja sistem pada pengujian satu tahap diukur menggunakan serangkaian metrik yang dirancang secara spesifik untuk menilai sistem agentik pada berbagai aspek. Mayoritas metrik yang digunakan, yaitu \textit{Disease Accuracy}, \textit{Trajectory Accuracy}, \textit{Agent Goal Accuracy}, dan \textit{Faithfulness}, merupakan metrik berbasis \textit{LLM-as-a-judge} yang bersifat \textit{self-explaining}.

\begin{enumerate}[label=\arabic*., listparindent=\parindent, labelindent=0pt]

\item \textit{Disease Accuracy}

\textit{Disease Accuracy} ($M_{da}$) merupakan metrik kustom yang menjadi tolok ukur utama dalam penelitian ini untuk mengukur kemampuan sistem dalam mengidentifikasi penyakit tanaman. Berbeda dengan metrik klasifikasi konvensional yang hanya menilai ketepatan label kelas secara biner, metrik ini dirancang untuk melakukan penilaian secara bernuansa sesuai dengan prinsip fitopatologi. Dalam identifikasi penyakit tanaman, identifikasi tidak selalu harus menghasilkan nama penyakit yang persis sama, namun juga dianggap berhasil apabila sistem mampu mengidentifikasi sumber gejala, misalnya jenis patogen atau kategori penyakit dengan benar meskipun nama spesifiknya sedikit berbeda.

Untuk memastikan penilaian yang objektif dan konsisten, metrik ini diimplementasikan menggunakan struktur \textit{Directed Acyclic Graph} (DAG) dari DeepEval. Penggunaan DAG bertujuan untuk membuat proses penilaian oleh \textit{LLM-as-a-judge} menjadi deterministik. Data yang diperlukan untuk pengukuran \textit{Disease Accuracy} disajikan pada Tabel \ref{tab:disease_accuracy_data}.

{
\setstretch{1.0}
\begin{longtblr}[
    caption = {Kebutuhan data untuk pengukuran \textit{disease accuracy}},
    label = {tab:disease_accuracy_data},
]{
    width = \linewidth,
    colspec = {X[0.8,l] X[1.7,l] X[3,l]},
    hlines, vlines,
    rows = {m},
    rowsep = 3pt,
    row{1} = {font=\bfseries, halign=c, valign=m},
    column{1} = {font=\bfseries},
    rowhead = 1,
}
    Kategori & Parameter & Deskripsi \\
    \SetCell[r=2]{l} Input & \textit{user\_text} & Teks pertanyaan atau keluhan yang diberikan pengguna. \\
     & \textit{image\_url} & Tautan citra tanaman yang diunggah pengguna. \\
    Output & \textit{final\_answer} & Jawaban akhir yang dihasilkan oleh agen. \\
    \SetCell[r=3]{l} Referensi & \textit{expected\_disease} & Nama penyakit yang sebenarnya (\textit{ground truth}) atau label "TANAMAN SEHAT". \\
     & \textit{pathogen\_type} & Jenis patogen penyebab penyakit (misal: Fungi, Bakteri, Virus). \\
     & \textit{reference\_answer} & Jawaban referensi yang lengkap dan benar sebagai pembanding. \\
\end{longtblr}
}

Logika penilaian dalam DAG disusun secara hierarkis, diawali dengan verifikasi identifikasi tanaman sehat. Tahap selanjutnya memeriksa akurasi identifikasi, mulai dari penyakit utama, kecocokan parsial, hingga jenis patogen. Apabila identifikasi gagal, penilaian beralih pada ketepatan penundaan jawaban (\textit{abstention}) akibat kurangnya bukti. Visualisasi alur ini ditunjukkan pada Gambar \ref{fig:disease_accuracy_dag}.

\noindent
    \begin{minipage}{\linewidth}
        \vspace{0.4cm}
        \captionsetup{type=figure}
        \centering
        \includegraphics[width=0.9\textwidth]{bab3/dagtree.png}
        \caption{Diagram logika penilaian DAG dari \textit{disease accuracy}}
        \label{fig:disease_accuracy_dag}
        \vspace{0.6cm}
    \end{minipage}

\item \textit{Trajectory Accuracy}

\textit{Trajectory Accuracy} ($M_{ta}$) adalah metrik yang dirancang untuk mengukur efisiensi dan ketepatan urutan pengambilan keputusan (\textit{trajectory}) oleh agen dalam menyelesaikan tugas. Pengukuran ini menggunakan pendekatan \textit{LLM-as-a-judge}, di mana LLM bertindak sebagai penilai kualitatif yang memvalidasi keputusan agen terhadap rubrik yang telah ditentukan. Pendekatan ini diterapkan karena mekanisme kerja sistem agentik yang kompleks, yang melibatkan mekanisme cadangan (\textit{fallback}) pada modul deteksi objek maupun pencarian informasi yang bergantung pada konteks situasi.

Metrik ini dirancang menggunakan kerangka kerja AgentEval dari LangSmith. Penilaian dilakukan dengan membandingkan trajektori aktual agen terhadap trajektori referensi menggunakan prinsip \textit{superset} dan kesesuaian semantik. Prinsip \textit{superset} mengizinkan agen untuk memanggil alat tambahan yang relevan, misalnya alat deteksi tambahan untuk validasi, selama seluruh alat kunci yang terdapat dalam referensi telah dipenuhi. Sementara itu, prinsip kesesuaian semantik memungkinkan evaluator untuk memvalidasi penggunaan alat yang berbeda namun memiliki fungsi yang setara, seperti penggunaan pencarian web (\textit{web\_search}) untuk menggantikan pencarian basis pengetahuan (\textit{knowledgebase\_search}) apabila informasi internal tidak memadai. Tabel \ref{tab:trajectory_data_requirements} menyajikan kebutuhan data untuk pengukuran \textit{Trajectory Accuracy}.

{
\setstretch{1.0}
\begin{longtblr}[
    caption = {Kebutuhan data untuk pengukuran \textit{trajectory accuracy}},
    label = {tab:trajectory_data_requirements},
]{
    width = \linewidth,
    colspec = {X[1,l] X[1.5,l] X[3,l]},
    hlines, vlines,
    rows = {m},
    rowsep = 3pt,
    row{1} = {font=\bfseries, halign=c, valign=m},
    column{1} = {font=\bfseries},
    rowhead = 1,
}
    Kategori & Parameter & Deskripsi \\
    \SetCell[r=2]{l} Input & \textit{user\_text} & Instruksi atau pertanyaan teks dari pengguna. \\
     & \textit{image\_url} & Tautan citra tanaman yang menjadi konteks visual analisis. \\
    Aktual & \textit{trace\_messages} & Riwayat lengkap langkah eksekusi (\textit{trace}) dan pemanggilan alat oleh agen. \\
    Referensi & \textit{reference\_tool\_calls} & Daftar alat yang diharapkan dipanggil untuk menyelesaikan tugas secara optimal (Standar Emas). \\
\end{longtblr}
}

Penilaian dilakukan dengan memberikan instruksi (\textit{prompt}) spesifik kepada LLM penilai yang mencakup prinsip desain agen, aturan semantik, dan rubrik penilaian. Model Gemini 2.5 Pro digunakan sebagai evaluator untuk memastikan kemampuan penalaran yang tinggi dalam menilai konteks trajektori. Rangkuman logika penilaian dan rubrik skor disajikan pada Tabel \ref{tab:trajectory_rubric}.

{
\setstretch{1.0}
\begin{longtblr}[
    caption = {Logika penilaian dan rubrik skor \textit{trajectory accuracy}},
    label = {tab:trajectory_rubric},
]{
    width = \linewidth,
    colspec = {X[1.5,l] X[1.2,l] X[3,l]},
    hlines, vlines,
    rows = {m},
    rowsep = 3pt,
    row{1} = {font=\bfseries, halign=c, valign=m},
    column{1} = {font=\bfseries},
    rowhead = 1,
}
    Komponen & Elemen & Deskripsi \\
    \SetCell[r=2]{l} Prinsip Desain & Alur Kerja & Identifikasi visual diprioritaskan, diikuti deteksi opsional, dan pencarian pengetahuan jika pengguna membutuhkan panduan pengobatan. \\
     & Fallback & Pengalihan ke pencarian jika kepercayaan visual rendah. \\
    \SetCell[r=2]{l} Aturan Semantik & Superset & Agen boleh memanggil alat tambahan (misal: deteksi ekstra) selama alat wajib terpenuhi. \\
     & Substitusi & \textit{knowledgebase\_search} dan \textit{web\_search} dianggap setara sebagai alat pencarian informasi. \\
    \SetCell[r=5]{l} Rubrik Skor & 0.0 & Alat kritis referensi tidak dipanggil sama sekali. \\
     & 0.2 & Kehilangan beberapa alat kunci penting. \\
     & 0.5 & Alat wajib lengkap, namun terdapat langkah tambahan yang tidak logis atau berlebihan. \\
     & 0.75 & \textit{Superset} valid dengan inefisiensi minor (misal: urutan kurang optimal). \\
     & 1.0 & \textit{Superset} sempurna; semua alat wajib ada, langkah tambahan beralasan kuat, dan alur logis. \\
\end{longtblr}
}

Skor akhir yang dihasilkan berupa nilai numerik diskrit (0.0, 0.2, 0.5, 0.75, atau 1.0) yang merepresentasikan kualitas trajektori agen dalam mencapai solusi yang efisien dan logis.

\item \textit{Agent Goal Accuracy}

\textit{Agent Goal Accuracy} ($M_{aga}$) adalah metrik kustom yang digunakan untuk mengukur eksekusi agen dan menentukan apakah tujuan tingkat tinggi (\textit{high-level goal}) telah tercapai. Berbeda dengan \textit{Trajectory Accuracy} yang lebih menitikberatkan pada validitas urutan pemanggilan alat, metrik ini mengukur pencapaian tujuan secara komprehensif dengan mempertimbangkan masukan pengguna, konteks pemanggilan alat, serta kualitas jawaban akhir. Penilaian dilakukan dengan membandingkan hasil aktual terhadap tujuan referensi yang merepresentasikan hasil ideal dari interaksi.

Metrik ini dirancang menggunakan kerangka kerja G-Eval yang diusulkan oleh \textcite{liuGEvalNLGEvaluation2023} melalui pustaka DeepEval. G-Eval digunakan untuk mengurangi subjektivitas dari evaluator tradisional dengan memanfaatkan LLM (Gemini 2.5 Pro) sebagai penilai yang mengikuti kriteria evaluasi terstruktur (\textit{LLM-as-a-judge}). Kebutuhan data untuk pengukuran \textit{Agent Goal Accuracy} dirinci pada Tabel \ref{tab:goal_accuracy_data}.

{
\setstretch{1.0}
\begin{longtblr}[
    caption = {Kebutuhan data untuk pengukuran \textit{agent goal accuracy}},
    label = {tab:goal_accuracy_data},
]{
    width = \linewidth,
    colspec = {X[1,l] X[1.5,l] X[3,l]},
    hlines, vlines,
    rows = {m},
    rowsep = 3pt,
    row{1} = {font=\bfseries, halign=c, valign=m},
    column{1} = {font=\bfseries},
    rowhead = 1,
}
    Kategori & Parameter & Deskripsi \\
    \SetCell[r=2]{l} Input & \textit{user\_text} & Instruksi atau pertanyaan teks dari pengguna. \\
     & \textit{image\_url} & Tautan citra masukan (\textit{multimodal}). \\
    \SetCell[r=2]{l} Aktual & \textit{final\_answer} & Jawaban akhir yang dihasilkan oleh agen. \\
     & \textit{tools\_called} & Daftar alat yang dipanggil agen beserta keluarannya. \\
    Referensi & \textit{reference\_goal} & Tujuan referensi yang diharapkan (misal: identifikasi, panduan, atau penolakan). \\
\end{longtblr}
}

Desain evaluator ini mencakup kriteria, langkah evaluasi, dan rubrik penilaian yang dirangkum pada Tabel \ref{tab:goal_evaluator_design}.

{
\setstretch{1.0}
\begin{longtblr}[
    caption = {Desain evaluator \textit{agent goal accuracy}},
    label = {tab:goal_evaluator_design},
]{
    width = \linewidth,
    colspec = {X[1,l] X[3.5,l]},
    hlines, vlines,
    rows = {m},
    rowsep = 3pt,
    row{1} = {font=\bfseries, halign=c, valign=m},
    column{1} = {font=\bfseries},
    rowhead = 1,
}
    Komponen & Rincian \\
    Kriteria & Evaluasi keberhasilan agen mencapai \textit{reference\_goal} berdasarkan output aktual dan penggunaan alat. Penilaian mencakup aspek identifikasi penyakit, panduan manajemen, analisis gejala, ekspresi ketidakpastian, dan penolakan permintaan di luar jangkauan. \\
    Langkah Evaluasi & \begin{minipage}{\linewidth}\begin{enumerate}[label=\arabic*., listparindent=\parindent, labelindent=0pt, nosep, leftmargin=*]
        \item Ekstrak tujuan referensi dari data yang diharapkan.
        \item Identifikasi tujuan utama (identifikasi, pengobatan, atau penolakan).
        \item Periksa alat yang dipanggil dan outputnya.
        \item Verifikasi observasi visual agen (jika gambar tersedia).
        \item Evaluasi kualitas dan kelengkapan pemenuhan tujuan dalam respon akhir.
    \end{enumerate}\end{minipage} \\
    Rubrik & \begin{minipage}{\linewidth}\begin{itemize}[nosep, leftmargin=*]
        \item 0--1 (Tidak Ada): Agen gagal menangani tujuan atau melenceng sepenuhnya.
        \item 2--4 (Minimal): Mencoba mencapai tujuan tetapi dengan celah signifikan atau penggunaan alat yang buruk.
        \item 5--8 (Parsial): Menangani tujuan utama tetapi melewatkan komponen kunci atau kurang mendalam.
        \item 9--10 (Penuh): Sepenuhnya memenuhi tujuan referensi dengan kualitas tinggi dan penggunaan alat yang tepat.
    \end{itemize}\end{minipage} \\
\end{longtblr}
}

Skor yang dihasilkan oleh evaluator G-Eval dinormalisasi ke dalam rentang 0.0 hingga 1.0, di mana nilai 1.0 menunjukkan pencapaian tujuan yang sempurna sesuai referensi.

\item \textit{Faithfulness}

\textit{Faithfulness} ($M_{fth}$) merupakan metrik kustom yang mengadopsi pendekatan \textit{LLM-as-a-judge} menggunakan kerangka kerja G-Eval dari pustaka DeepEval. Metrik ini berfungsi untuk mengukur integritas faktual dari respons agen, memastikan bahwa jawaban akhir (\textit{final answer}) yang dihasilkan selaras dengan konteks informasi yang diperoleh dari seluruh alat (\textit{retrieved context}). Pengukuran ini krusial untuk mendeteksi halusinasi, yaitu kondisi di mana model menghasilkan informasi yang tampak meyakinkan namun tidak didukung oleh bukti data internal maupun eksternal yang diakses.

Berbeda dengan implementasi standar, metrik ini dirancang khusus untuk mengakomodasi sifat \textit{multimodal} dan pengetahuan domain agen. Dalam desain ini, observasi visual yang dilakukan agen terhadap citra tanaman serta pemberian saran perawatan umum yang bersifat standar (\textit{best practices}) diakui sebagai informasi valid yang tidak memerlukan sitasi dari alat pencarian, selama masih dalam batas kewajaran domain. Tabel \ref{tab:faithfulness_data_requirements} menyajikan komponen data yang diperlukan dalam proses pengukuran.

{
\setstretch{1.0}
\begin{longtblr}[
    caption = {Kebutuhan data untuk pengukuran \textit{faithfulness}},
    label = {tab:faithfulness_data_requirements},
]{
    width = \linewidth,
    colspec = {X[1,l] X[1.5,l] X[3,l]},
    hlines, vlines,
    rows = {m},
    rowsep = 3pt,
    row{1} = {font=\bfseries, halign=c, valign=m},
    column{1} = {font=\bfseries},
    rowhead = 1,
}
    Kategori & Parameter & Deskripsi \\
    \SetCell[r=2]{l} Input & \textit{user\_text} & Pertanyaan atau instruksi teks dari pengguna. \\
     & \textit{image\_url} & Tautan citra tanaman sebagai dasar validasi observasi visual agen. \\
    \SetCell[r=2]{l} Aktual & \textit{final\_answer} & Jawaban akhir yang dihasilkan oleh agen. \\
     & \textit{retrieval\_context} & Kumpulan potongan informasi teks yang diekstraksi dari keluaran seluruh alat (pencarian, deteksi, dan klasifikasi). \\
\end{longtblr}
}

Desain evaluator \textit{Faithfulness} mencakup kriteria penilaian yang spesifik terhadap perilaku agen serta rubrik penilaian bertingkat. Rincian desain ini dirangkum pada Tabel \ref{tab:faithfulness_evaluator_design}.

{
\setstretch{1.0}
\begin{longtblr}[
    caption = {Desain evaluator \textit{faithfulness}},
    label = {tab:faithfulness_evaluator_design},
]{
    width = \linewidth,
    colspec = {X[1,l] X[3.5,l]},
    hlines, vlines,
    rows = {m},
    rowsep = 3pt,
    row{1} = {font=\bfseries, halign=c, valign=m},
    column{1} = {font=\bfseries},
    rowhead = 1,
}
    Komponen & Rincian \\
    Kriteria & Evaluasi kesetiaan (\textit{faithfulness}) respons terhadap bukti dengan mempertimbangkan kapabilitas agen:
    \begin{minipage}{\linewidth}\begin{enumerate}[label=\arabic*., listparindent=\parindent, labelindent=0pt, nosep, leftmargin=*]
        \item Observasi Visual: Agen HARUS melakukan observasi dari gambar (valid tanpa pencarian).
        \item Saran Umum: Agen DAPAT memberikan praktik terbaik umum (sanitasi, isolasi) tanpa pencarian.
        \item Klaim Spesifik: Identifikasi penyakit, nama patogen, dan pengobatan kimiawi HARUS didasarkan pada \textit{retrieval\_context}.
        \item Ketidakpastian: Agen HARUS menyatakan ketidakpastian jika bukti terbatas.
    \end{enumerate}\end{minipage} \\
    Langkah Evaluasi & \begin{minipage}{\linewidth}\begin{enumerate}[label=\arabic*., listparindent=\parindent, labelindent=0pt, nosep, leftmargin=*]
        \item Identifikasi observasi visual dan rekomendasi umum dalam output.
        \item Identifikasi semua klaim spesifik penyakit (nama, biologi, kimia).
        \item Verifikasi apakah observasi visual masuk akal terhadap citra.
        \item Periksa apakah setiap klaim spesifik didukung oleh konteks pencarian.
        \item Hitung proporsi klaim yang tidak didukung atau difabrikasi.
    \end{enumerate}\end{minipage} \\
    Rubrik & \begin{minipage}{\linewidth}\begin{itemize}[nosep, leftmargin=*]
        \item 0.0--0.2 (Tidak Setia): Informasi didominasi fabrikasi; rekomendasi spesifik dikarang tanpa dasar.
        \item 0.3--0.4 (Sebagian Besar Tidak Setia): Banyak klaim spesifik tidak didukung bukti.
        \item 0.5--0.6 (Setia Sebagian): Beberapa klaim kurang dukungan atau spekulasi disajikan sebagai fakta.
        \item 0.7--0.8 (Sebagian Besar Setia): Klaim inti didasarkan bukti, hanya detail minor yang terlewat.
        \item 0.9--1.0 (Sangat Setia): Seluruh klaim spesifik didukung bukti; observasi visual dan saran umum tepat.
    \end{itemize}\end{minipage} \\
\end{longtblr}
}

Skor akhir yang dihasilkan oleh evaluator direpresentasikan dalam nilai kontinu antara 0.0 hingga 1.0, di mana nilai yang lebih tinggi mengindikasikan tingkat kepatuhan fakta yang lebih baik terhadap sumber data.

\item \textit{Context Relevance}

\textit{Context Relevance} ($M_{cr}$) merupakan metrik yang diadopsi dari kerangka kerja Ragas untuk mengukur seberapa relevan potongan informasi (\textit{retrieved contexts}) yang diambil oleh sistem terhadap masukan pengguna (\textit{user input}). Pengukuran ini dilakukan menggunakan mekanisme \textit{LLM-as-a-judge} yang melibatkan dua pemanggilan \textit{prompt} penilaian secara independen untuk mengurangi bias model. Setiap \textit{prompt} menginstruksikan model evaluator untuk memberikan penilaian relevansi pada skala diskret 0 hingga 2. Nilai 0 mengindikasikan bahwa konteks yang diambil sama sekali tidak relevan dengan kueri, nilai 1 menunjukkan relevansi sebagian, dan nilai 2 menandakan bahwa konteks sepenuhnya relevan. Skor dari kedua penilaian tersebut kemudian dinormalisasi ke dalam rentang [0,1] dan dirata-ratakan untuk menghasilkan skor akhir. Skor yang lebih tinggi merepresentasikan keselarasan yang lebih baik antara konteks informasi yang ditemukan dengan kebutuhan pengguna.

Metrik ini berbeda dengan pendekatan \textit{Context Relevance} awal yang diperkenalkan pada makalah Ragas oleh \textcite{esRagasAutomatedEvaluation2025}, yang mendefinisikan relevansi berdasarkan rasio ekstraksi tingkat kalimat (jumlah kalimat relevan dibagi total kalimat). Implementasi pada penelitian ini menggunakan pendekatan penilaian diskret (\textit{discrete rating}) yang dinilai lebih \textit{robust} dan efisien. Pendekatan ini mengurangi kerentanan terhadap kesalahan pemisahan kalimat (\textit{sentence boundary errors}) yang sering terjadi pada teks teknis, serta lebih hemat biaya komputasi (\textit{overhead}) sambil tetap mempertahankan tujuan inti pengukuran. Selain itu, implementasi metrik ini telah disesuaikan untuk mendukung input \textit{multimodal}, di mana konteks visual dari citra turut diperhitungkan dalam formulasi pertanyaan pengguna yang disajikan ke evaluator.

Sebagai ilustrasi perhitungan, misalkan pengguna bertanya "Apa obat untuk penyakit bercak daun ini?" dan sistem mengambil dua konteks: (1) "Bercak daun Cercospora dapat dikendalikan dengan fungisida berbahan aktif mankozeb" dan (2) "Sanitasi lahan penting untuk mencegah penyebaran". Karena kedua konteks tersebut secara kolektif menjawab pertanyaan pengguna (memberikan solusi obat dan pencegahan), kedua \textit{prompt} evaluator akan memberikan skor 2 (sepenuhnya relevan). Skor ini kemudian dinormalisasi ($2/2 = 1.0$) dan dirata-ratakan, menghasilkan skor akhir \textit{Context Relevance} sebesar 1.0. Kebutuhan data untuk pengukuran \textit{Context Relevance} dirangkum pada Tabel \ref{tab:context_relevance_data}.

{
\setstretch{1.0}
\begin{longtblr}[
    caption = {Kebutuhan data untuk pengukuran \textit{context relevance}},
    label = {tab:context_relevance_data},
]{
    width = \linewidth,
    colspec = {X[1,l] X[1.5,l] X[3,l]},
    hlines, vlines,
    rows = {m},
    rowsep = 3pt,
    row{1} = {font=\bfseries, halign=c, valign=m},
    column{1} = {font=\bfseries},
    rowhead = 1,
}
    Kategori & Parameter & Deskripsi \\
    \SetCell[r=2]{l} Input & \textit{user\_text} & Teks pertanyaan atau instruksi dari pengguna. \\
     & \textit{image\_url} & Tautan citra tanaman. Jika tersedia, indikator visual ditambahkan ke \textit{prompt} evaluator untuk konteks. \\
    Aktual & \textit{retrieved\_contexts} & Daftar potongan teks yang diekstraksi dari riwayat pemanggilan alat (\textit{trace}), mencakup hasil pencarian dokumen maupun deteksi. \\
\end{longtblr}
}

\end{enumerate}
\end{enumerate}

\subsubsection{Pengujian Multi-Tahap}
Pengujian multi-tahap dalam penelitian ini difokuskan pada eksperimen batasan domain (\textit{Domain Boundary}). Eksperimen ini bertujuan untuk menguji ketahanan sistem dalam menjaga konsistensi peran dan relevansi topik selama interaksi percakapan yang berlangsung dalam beberapa giliran (\textit{multi-turn}). Pengujian ini diorkestrasi menggunakan fitur simulasi interaksi dari platform LangSmith \parencite{langchain_multiturn_simulation}, di mana agen dihadapkan dengan "pengguna simulasi" (\textit{simulated user}) yang dijalankan oleh LLM lain dengan instruksi persona tertentu. Pendekatan ini memungkinkan evaluasi yang dinamis terhadap respons agen saat menghadapi berbagai skenario percakapan yang berkembang.

\begin{enumerate}[label=\alph*., listparindent=\parindent, labelindent=0pt]
    \item Persiapan Dataset Pengujian Multi-Tahap

    Persiapan dataset pengujian dilakukan dengan menyusun dataset yang berisi 10 contoh skenario interaksi. Setiap skenario dilengkapi dengan pesan awal pengguna dan instruksi sistem untuk pengguna simulasi guna membentuk persona yang spesifik. Variasi persona yang digunakan dalam pengujian ini dikelompokkan sebagai berikut:

\begin{enumerate}[label=\alph*., listparindent=\parindent, labelindent=0pt]
    \item \textit{On-topic Personas}: Persona pengguna yang mengajukan pertanyaan relevan mengenai kesehatan tanaman, baik disertai bukti visual maupun teks saja. Contoh persona ini meliputi Petani Tomat, Penghobi Mawar, dan Petani Cabai.
    \item \textit{Off-topic Personas}: Persona pengguna yang secara eksplisit meminta bantuan di luar domain keahlian agen. Contohnya mencakup pemilik mobil yang meminta saran mekanis, investor yang menanyakan harga saham, dan pertanyaan medis terkait kesehatan manusia.
    \item \textit{Edge Cases}: Skenario batas yang menguji kemampuan diskriminasi visual agen, seperti pengguna yang mengirimkan foto tanaman artifisial atau objek hewan, misalnya kucing, untuk melihat apakah agen dapat membedakan objek non-tanaman.
    \item \textit{Context Switchers}: Persona yang dirancang untuk menguji fleksibilitas agen dalam menangani perpindahan topik. Kategori ini mencakup perpindahan dari topik tanaman ke topik lain (\textit{On-to-Off}) untuk menguji ketegasan agen, serta perpindahan dari topik umum ke topik tanaman (\textit{Off-to-On}) untuk menguji responsivitas agen terhadap pertanyaan valid setelah penolakan awal.
\end{enumerate}

\item Metrik Pengukuran Pengujian Multi-Tahap

Kinerja agen pada pengujian multi-tahap diukur menggunakan metrik kustom bernama \textit{Domain Adherence} ($M_{dad}$). Metrik ini merupakan implementasi dari pendekatan \textit{LLM-as-a-judge} yang dibangun menggunakan pustaka OpenEval dari LangSmith dan bersifat \textit{self-explaining}, di mana evaluator tidak hanya memberikan skor penilaian tetapi juga menyertakan alasan kualitatif yang mendasari skor tersebut.

Evaluator bekerja dengan menganalisis riwayat percakapan lengkap (\textit{trajectory}) antara agen dan pengguna simulasi. Penilaian dilakukan berdasarkan instruksi (\textit{prompt}) yang menginstruksikan model juri untuk memberikan skor tinggi jika agen berhasil menjawab pertanyaan terkait tanaman dan menolak permintaan di luar domain secara sopan. Sebaliknya, skor rendah diberikan jika agen gagal mempertahankan batasan domainnya. \textit{Prompt} yang digunakan oleh evaluator dalam metrik ini disajikan pada Gambar \ref{fig:prompt_domain_adherence}.

\noindent
\begin{minipage}{\linewidth}
    \vspace{0.4cm}
    \captionsetup{type=figure}
    \centering
    \includegraphics[width=\textwidth]{bab3/promptdomainadherence.png}
    \caption{\textit{Prompt} penilaian untuk metrik \textit{Domain Adherence}}
    \label{fig:prompt_domain_adherence}
    \vspace{0.6cm}
\end{minipage}

\end{enumerate}


\section{Evaluasi}
Tahap evaluasi dilakukan untuk menganalisis hasil pengujian secara mendalam guna mengukur kualitas dan kinerja sistem yang telah dibangun. Analisis ini mencakup penilaian kualitatif untuk memvalidasi fungsionalitas fitur, serta penilaian kuantitatif untuk menentukan konfigurasi optimal dan mengukur efektivitas sistem berdasarkan metrik-metrik yang telah ditetapkan.

\subsection{Modul Deteksi Objek}
Evaluasi integrasi alat di modul ini pada sistem hanya dilakukan secara kualitatif berdasarkan hasil pengujian \textit{black box}. Sementara itu, pemilihan varian YOLOv11 terbaik pada alat deteksi \textit{closed-set} daun secara khusus dilakukan dengan menganalisis hubungan antara performa deteksi dan efisiensi latensi menggunakan \textit{Pareto frontier}. Pendekatan ini digunakan untuk mengidentifikasi sekumpulan model yang tidak dapat ditingkatkan pada salah satu aspek tanpa mengorbankan aspek lainnya, sehingga menghasilkan himpunan solusi optimal dalam ruang dua dimensi mAP dan latensi \parencite{shakyaSmallObjectDetectionEdge2026}. Model-model yang berada pada kurva \textit{Pareto frontier} merepresentasikan pilihan terbaik untuk berbagai preferensi antara akurasi dan kecepatan. Varian yang dipilih untuk diintegrasikan ke dalam sistem agentik adalah model yang berada pada kurva tersebut sekaligus menawarkan latensi paling minimal dengan tetap mempertahankan performa deteksi yang memadai. Visualisasi contoh grafik \textit{Pareto frontier} disajikan pada Gambar \ref{fig:pareto_curve}.
mempertahankan performa deteksi yang memadai. Visualisasi contoh grafi
\begin{figure}[H]
    \centering
    \includegraphics[width=0.9\textwidth]{bab3/contohparetofrontier.png}
    \caption[Contoh grafik \textit{Pareto Frontier} mAP berbanding latensi]{Contoh grafik \textit{Pareto Frontier} mAP berbanding latensi \parencite{shakyaSmallObjectDetectionEdge2026}}
    \label{fig:pareto_curve}
\end{figure}

\subsection{Modul Pengambilan Informasi}
Evaluasi pada modul pengambilan informasi dibedakan berdasarkan cara pengujiannya. Alat pengambilan citra-teks dan alat pencarian internet yang diuji secara \textit{black box} dianalisis secara kualitatif untuk memvalidasi fungsionalitasnya. Skenario pengujian beserta hasil yang diharapkan dari kedua alat tersebut dirangkum pada Tabel \ref{tab:black_box_modul_pengambilan}.

{
\setstretch{1.0}
\begin{longtblr}[
    caption = {Pengujian \textit{black box} modul pengambilan informasi},
    entry = {Pengujian \textit{black box} Modul Pengambilan Informasi},
    label = {tab:black_box_modul_pengambilan},
]{
    width = \linewidth,
    colspec = {X[0.6,l] X[1.2,l] X[2,l] X[2.5,l]},
    hlines, vlines,
    rows = {m},
    rowhead = 1,
    row{1} = {font=\bfseries, halign=c, valign=m},
    rowsep = 3pt,
}
    ID & Alat & Skenario Pengujian & Hasil yang Diharapkan \\
    PI-BB1 & Pengambilan Citra-Teks & \textit{Query} teks diberikan sebagai masukan untuk teknik \textit{text-to-text}. & Sistem mengembalikan dokumen teks yang relevan secara semantik dari basis data vektor. \\
    PI-BB2 & Pengambilan Citra-Teks & \textit{Query} teks diberikan sebagai masukan untuk teknik \textit{text-to-image}. & Sistem mengembalikan citra penyakit tanaman yang relevan dari basis data vektor. \\
    PI-BB3 & Pengambilan Citra-Teks & Citra penyakit tanaman diberikan sebagai masukan untuk teknik \textit{image-to-image}. & Sistem mengembalikan citra lain yang memiliki kemiripan visual tinggi dari basis data vektor. \\
    PI-BB4 & Pengambilan Citra-Teks & Citra penyakit tanaman diberikan sebagai masukan untuk teknik \textit{image-to-text}. & Sistem mengembalikan dokumen teks deskriptif yang relevan secara semantik dengan konten citra. \\
    PI-BB5 & Pengambilan Citra-Teks & \textit{Query} diberikan disertai filter nama tanaman tertentu. & Sistem hanya mengembalikan hasil yang berasal dari tanaman yang sesuai dengan nilai filter yang ditentukan. \\
    PI-BB6 & Pengambilan Citra-Teks & Citra beserta koordinat \textit{bounding box} hasil deteksi objek diberikan sebagai masukan. & Sistem melakukan pengambilan berbasis wilayah (\textit{region-based}) pada area yang ditentukan dan mengembalikan hasil yang lebih terfokus pada bagian tanaman tersebut. \\
    PI-BB7 & Pencarian Internet & \textit{Query} teks diberikan sebagai masukan untuk pencarian web. & Sistem berhasil mengembalikan daftar hasil pencarian dari internet yang relevan dengan \textit{query}. \\
\end{longtblr}
}

Sementara itu, alat pengambilan basis data pengetahuan dianalisis secara kuantitatif. Analisis dilakukan pada metrik MRR, NDCG@5, dan Recall@5 yang diperoleh dari pengujian. Analisis kuantitatif ini bertujuan untuk mengidentifikasi seberapa jauh perbedaan kinerja antarkonfigurasi serta menentukan konfigurasi yang paling optimal secara terukur. Evaluasi dilakukan dalam dua tahap berurutan. Pertama, berbagai kombinasi teknik \textit{hybrid search} dibandingkan untuk menentukan konfigurasi dasar (\textit{base retriever}) terbaik, kemudian dampak penambahan komponen \textit{reranker} dievaluasi terhadap konfigurasi terbaik tersebut.

\subsection{Antarmuka Sistem Agentik}
Evaluasi antarmuka sistem agentik dilakukan secara kualitatif berdasarkan hasil pengujian \textit{black box} yang telah dilaksanakan pada tahap sebelumnya. Analisis ini bertujuan untuk memvalidasi bahwa seluruh fungsionalitas yang dirancang pada Gambar \ref{fig:use_case_diagram} telah diimplementasikan dengan benar dan berfungsi sesuai dengan spesifikasi yang ditetapkan. Setiap skenario pengujian yang dirangkum pada Tabel \ref{tab:black_box_testing_antarmuka} dievaluasi berdasarkan kesesuaian keluaran aktual sistem dengan hasil yang diharapkan, sehingga dapat ditentukan apakah antarmuka sistem agentik telah memenuhi persyaratan fungsional yang telah dirancang.

\subsection{Sistem Agentik}
Hasil pengujian sistem agentik dianalisis secara kuantitatif berdasarkan skor yang dihasilkan oleh masing-masing metrik pengukuran yang telah ditetapkan. Analisis mencakup pemeriksaan kasus pengujian yang berhasil dan gagal dengan mengacu pada ambang skor tiap metrik, serta identifikasi pola kesalahan yang muncul secara konsisten. Pada pengujian satu tahap, perbandingan dilakukan antara dataset ID dan OOD untuk mengukur kemampuan generalisasi sistem. Sementara itu, pada pengujian multi-tahap, analisis diarahkan pada konsistensi agen dalam mempertahankan batasan domain selama percakapan berlangsung.

\section{Penyampaian Hasil}
Hasil penelitian yang telah dilaksanakan dianalisis dan dibahas dalam bentuk laporan tertulis berupa skripsi. Laporan ini mencakup perancangan arsitektur sistem agentik, implementasi teknis, serta hasil evaluasi kinerja sistem. Pembahasan dalam laporan meliputi integrasi model deteksi objek menggunakan YOLOv11 dan OWLv2, model pengambilan citra-teks \textit{multimodal} berbasis SCOLD, serta pengembangan agen cerdas dengan pola ReAct yang memanfaatkan mekanisme RAG. Laporan ini diharapkan dapat memberikan kontribusi terhadap pengembangan teknologi identifikasi penyakit tanaman yang lebih adaptif dan fleksibel. Selain itu, penelitian ini juga diharapkan dapat menjadi rujukan dalam penerapan teknologi sistem agentik dan AI \textit{multimodal} di bidang pertanian untuk membantu pengelolaan kesehatan tanaman.